\documentclass[11pt,a4paper,twoside,openright]{report}

%  use packages
% --------------
\usepackage[utf8]{inputenc}
\usepackage{geometry,amsmath, amsthm, amssymb}   % , units, dsfont, caption
% enumerate: \begin{enumerate}[{Example} I)] ... Roman numbering
\usepackage{graphicx,color,tabularx,float,subcaption,algorithmicx,algpseudocode,mathtools,textcomp,multirow,emptypage}
\usepackage[super]{nth}
\usepackage[table]{xcolor}
\usepackage[titletoc]{appendix}
\usepackage{nomencl}
\makenomenclature

\usepackage[pdftex,
						pdftitle={Side-Channel Attack Analysis of AES White-Box Schemes, Jakub Klemsa},
						pdfauthor={Jakub Klemsa},
						pdfsubject={},
						pdfkeywords={White-Box AES, Side-Channel Attack},
						pdfcreator=pdflatex]{hyperref}

% fix left/right margins
\let\tmp\oddsidemargin
\let\oddsidemargin\evensidemargin
\let\evensidemargin\tmp
\reversemarginpar

\usepackage{fancyhdr}
	\pagestyle{fancy}
	\fancyhf{}
	\fancyhead[L]{\slshape \rightmark}   % \leftmark .. kapitola
	\fancyhead[R]{\thepage}

%  customize
% -----------
\let\tmpitem\item
\renewcommand{\item}{\setlength{\itemsep}{0pt}\tmpitem}   % míň roztahovačný seznamy
\renewcommand{\labelitemi}{$\circ$}

\let\epsilon\varepsilon
\let\phi\varphi

\setcounter{tocdepth}{1}
%~ \setcounter{secnumdepth}{4}

\renewcommand{\bibname}{References}
\renewcommand{\contentsname}{Table of Contents}
\renewcommand{\nomname}{Notations and Symbols}

%  theorems etc.
% ---------------
\theoremstyle{definition}
\newtheorem{thm}{Theorem}[chapter]
\newtheorem{lemma}[thm]{Lemma}
\newtheorem{thesis}[thm]{Thesis}
\newtheorem{conj}[thm]{Conjecture}
\newtheorem{cor}[thm]{Corollary}
\newtheorem{defn}[thm]{Definition}
\newtheorem{example}[thm]{Example}
\newtheorem{game}[thm]{Game}
\newtheorem{princ}[thm]{Principle}
\newtheorem{alg}{Algorithm}[chapter]

\theoremstyle{remark} 
\newtheorem{remark}[thm]{Remark}
\newtheorem{note}[thm]{Note}
\newtheorem{notion}[thm]{Notion}

%  my commands
% -------------
\renewcommand{\iff}{\Leftrightarrow}
\newcommand{\then}{\Rightarrow}
\newcommand{\rarr}{\rightarrow}
\newcommand{\larr}{\leftarrow}
\newcommand{\xor}{\oplus}
\newcommand{\unirand}{\xleftarrow{R}}
%~ \DeclareMathOperator*{\Prob}{\mathbb{Pr}}   % operátor* bude umět používat \limit :) ... \Pr už existuje! a \mathbb neumí malý písmena
\DeclareMathOperator*{\Adv}{Adv}
\DeclareMathOperator*{\AdvPRP}{Adv_{PRP}}
\DeclareMathOperator*{\AdvSS}{Adv_{SS}}
\DeclareMathOperator*{\AdvPRPSS}{Adv_{PRPSS}}
\DeclareMathOperator*{\bin}{bin}
\DeclareMathOperator*{\Corr}{Corr}
\DeclareMathOperator*{\Mean}{Mean}

\renewcommand{\P}{\mathbb{P}}
\newcommand{\GF}{\mathrm{GF}}
\newcommand{\N}{\mathbb{N}}
\newcommand{\Z}{\mathbb{Z}}
\newcommand{\R}{\mathbb{R}}
\newcommand{\C}{\mathbb{C}}
\newcommand{\Pclass}{\mathsf{P}}
\newcommand{\NPclass}{\mathsf{NP}}

\newcommand{\SubBytes}{\texttt{SubBytes}}
\newcommand{\ShiftRows}{\texttt{ShiftRows}}
\newcommand{\MixColumns}{\texttt{MixColumns}}
\newcommand{\AddRoundKey}{\texttt{AddRoundKey}}
\newcommand{\ExpandedKey}{\texttt{ExpandedKey}}
\newcommand{\KeySchedule}{\texttt{KeySchedule}}

\newcommand{\TBox}{\texttt{TBox}}
\newcommand{\MC}{\texttt{MC}}
\newcommand{\MB}{\texttt{MB}}
\newcommand{\IMB}{\texttt{IMB}}
\newcommand{\Enc}{\texttt{Enc}}

\newcommand{\nicefrac}[2]{{}^{#1}\!/_{#2}}
\newcommand{\atob}[2]{#1 , \ldots , #2}
\newcommand{\link}[1]{\href{http://#1}{\tt #1}}

\newcommand{\cg}[1]{\cellcolor{gray!#1}}

%%%%%%%%%%%%%%%%%%%%%%%%%%%%%%%%%%%%%%%%%%%%%%%%%%%%%%%%%%%%%%%%%%%%%%%%
%%%   k vyplnění ZAČÁTEK
%%%%%%%%%%%%%%%%%%%%%%%%%%%%%%%%%%%%%%%%%%%%%%%%%%%%%%%%%%%%%%%%%%%%%%%%

\newcommand{\cvut}{ČESKÉ VYSOKÉ UČENÍ TECHNICKÉ V~PRAZE}
\newcommand{\fjfi}{Fakulta jaderná a fyzikálně inženýrská}
\newcommand{\km}{Katedra matematiky}
\newcommand{\obor}{Matematická informatika}
\newcommand{\TYPPRACE}{DIPLOMOVÁ PRÁCE}
\newcommand{\TypPrace}{Diplomová práce}
\newcommand{\mePrace}{mé diplomové práce}
\newcommand{\kohoCoPraci}{ji}

\newcommand{\nazevcz}{Analýza AES white-box schémat pomocí útoku postranním kanálem}
\newcommand{\nazeven}{Side-Channel Attack Analysis of AES White-Box Schemes}
\newcommand{\autor}{Jakub Klemsa}
\newcommand{\vedouci}{prof. RNDr. Václav Matyáš, M.Sc., Ph.D.}
\newcommand{\vedoucimu}{prof. RNDr. Václavu Matyáši, M.Sc., Ph.D.}
\newcommand{\pracovisteVed}{Fakulta informatiky, Masarykova Univerzita}
\newcommand{\konzultant}{Mgr. Dušan Klinec}
\newcommand{\rok}{2016}
\newcommand{\akadRok}{2015/2016}
\newcommand{\misto}{Praha}
\newcommand{\VMiste}{V Praze}

\newcommand{\klicova}{Whitebox AES, útok postranním kanálem, paměťová stopa.}
\newcommand{\keyword}{White-Box AES, Side-Channel Attack, Memory Trace.}
\newcommand{\abstrCZ}{Samotný koncept whitebox kryptografie je relativně nový, přesto už známe některé silné teoretické výsledky. Z praktického pohledu se ale jednalo a jedná spíše o hru kočky s myší -- návrhy se neopírají o žádný známý těžký problém a tak bývá spíše otázkou času, jak dlouho vydrží vzdorovat snaze kryptanalytiků.

Velmi zajímavý a univerzální přístup \cite{bos2015differential} byl představen v preprintu na \link{iacr.org}. Cílem této práce je naprogramovat potřebné nástroje a útok reprodukovat, předně na implementaci z diplomové práce Dušana Klince \cite{klinec2013white}. Dále budeme zjišťovat, kudy informace uniká a jak tomu zabránit.}
\newcommand{\abstrEN}{Although the concept of White-Box Cryptography is quite recent, we already know some strong theoretical results. But from the practical point of view it rather reminds a cat-and-mouse game -- proposed schemes are not supported by any known computationally hard problem so it is usually a matter of time until it gets broken by cryptanalytic effort.

A very interesting and universal approach has been recently introduced in a paper \cite{bos2015differential} at \link{iacr.org}'s preprint. The goal of this thesis is to write required tools and reproduce the attack, especially on an implementation from Dušan Klinec's diploma thesis. Then we will identify where the point of leakage is and how to prevent from leaking the key.}

%%%%%%%%%%%%%%%%%%%%%%%%%%%%%%%%%%%%%%%%%%%%%%%%%%%%%%%%%%%%%%%%%%%%%%%%
%%%   k vyplnění KONEC
%%%%%%%%%%%%%%%%%%%%%%%%%%%%%%%%%%%%%%%%%%%%%%%%%%%%%%%%%%%%%%%%%%%%%%%%


\begin{document}


\pagenumbering{roman}
\begin{titlepage}
	
	%%% 1. úvodní strana
	\thispagestyle{empty}
	\begin{center}
		{\Large \cvut \\[12pt] \fjfi \\[260pt]}
		{\Huge \TYPPRACE}
	\end{center}
	\vfill
	{
		\Large \misto, \rok \hfill \autor
	}
	\cleardoublepage
	
	
	%%% 2. úvodní strana
	\thispagestyle{empty}
	\begin{center}
		{\Large \cvut \\[10pt] \fjfi \\[10pt] \km \\[40pt]}
		\includegraphics[height=75pt]{lev.pdf} \\[100pt]
		
		{\LARGE \TYPPRACE \\[60pt]}
		
		{\Large\bf \nazevcz \\[30pt]}
		
		{\Large\bf \nazeven}
	\end{center}
	\vfill
	{
		\Large
		\begin{tabular}{ll}
		Vypracoval: & \autor\\[3pt]
		Školitel: & \vedouci\\[3pt]
		Akademický rok: & \akadRok
		\end{tabular}
	}
	\cleardoublepage
	
	
	%%% sem přijde ofic. zadání
	\thispagestyle{empty}
	\Large
	Na toto místo přijde svázat \textbf{zadání \mePrace}!
	
	\vspace{4mm}
	V~jednom z~výtisků musí být \textbf{originál} zadání, v~ostatních kopie.
	\normalsize
	\cleardoublepage
	
	
	%%% prohlášení
	\thispagestyle{empty}
	~
	\vfill
	\noindent\textbf{Čestné prohlášení}
	\vspace{0.5cm}
	
	\noindent
	Prohlašuji na tomto místě, že jsem předloženou práci vypracoval samostatně a že jsem uvedl veškerou použitou literaturu.
	\vspace{1.5cm}
	
	\noindent
	\vspace{5mm}\VMiste ~dne \today\hfill
		\begin{tabular}{c}
		. . . . . . . . . . . . . . . . . . . . . . .\\
		\autor
		\end{tabular}
	\cleardoublepage
	
	
	%%% poděkování
	\thispagestyle{empty}
	~
	\vfill
	\noindent\textbf{Poděkování}
	\vspace{0.5cm}
	
	\noindent
	Děkuji \vedoucimu, za vedení \mePrace ~a za podnětné návrhy, které \kohoCoPraci ~obohatily.
	
	\begin{flushright}
	\autor
	\end{flushright}
	\cleardoublepage
	
	
	%%% abstrakt atp.
	\thispagestyle{empty}
	\begin{tabularx}{\textwidth}{lX}
		{\em Název práce:} & \bf \nazevcz \\[4mm]
		{\em Autor:} & \autor \\[4mm]
		{\em Obor:} & \obor \\[4mm]
		{\em Druh práce:} & \TypPrace \\[4mm]
		{\em Vedoucí práce:} & \vedouci, \pracovisteVed \\[4mm]
		{\em Konzultant:} & \konzultant \\[4mm]
		{\em Abstrakt:} & \abstrCZ \\[4mm]
		{\em Klíčová slova:} & \klicova \\[20mm]
		
		{\em Title:} & \bf \nazeven \\[4mm]
		{\em Author:} & \autor \\[4mm]
		{\em Abstract:} & \abstrEN \\[4mm]
		{\em Keywords:} & \keyword
	\end{tabularx}
	\cleardoublepage
	
	
	%%% obsah
	\thispagestyle{empty}
	\tableofcontents
	\cleardoublepage
	
	
	%%% nomenklatura
	\thispagestyle{empty}
	\printnomenclature
	\cleardoublepage
	
\end{titlepage}

\pagenumbering{arabic}


%%%%%%%%%%%%%%%%%%%%%%%%%%%%%%%%%%%%%%%%%%%%%%%%%%%%%%%%%%%%%%%%%%%%%%%%
%%%   Z A Č Á T E K   P R Á C E
%%%%%%%%%%%%%%%%%%%%%%%%%%%%%%%%%%%%%%%%%%%%%%%%%%%%%%%%%%%%%%%%%%%%%%%%
		
		% FINAL CHECK:
		% 	abstrakty a keywords
		% 	V Praze dne \today -> český datum
		% 	\section{Capital Letters}
		% 	section \ref{...} -> Section \ref{...}
		% 	\mod vs \pmod
		% 	bib: title={{Bla Bla ... 2x závorky aby byly velký písmena}}
		% 	$SubBytes$ -> $\SubBytes$ ...
		% 	tables -> lookup tables
		% 	concatenation -- spelling
		% 	you vs. WE
		% 	an $8 ...
		% 	compose vs. concat !!!
		% 	E(k,m) instead of E(m,k)
		% 	adversary -> she
		% 	referred to as (-the?)
		% 	start -> begin with
		% 	závorky \left(\right)
		% 	setřídit nomenclaturu: \nomenclature[01]{$|S|$}{Cardinality of the set $S$}
		% 	| -> \mid (kde se hodí)
		% 	sofar -> so far
		% 	notion vs. notation
		% 	závorky \bigl( \bigr)
		% 	adversary vs. attacker
		% 	see vs. cf. (compare)
		% 	for this reason comma?
		% 	\Pr D E l k_a ??
		% 	e.g. -> e.g.\ 
		% 	všechno dát jako podmnožiny \{0,1\}^* !!!
		% 	přečíst si kydy z turku/sec eng a případně zařadit "historky z natáčení"
		% 	color box -> color-box
		% 	whitebox -> white-box
		% 	odkaz na section mimo chapter -> přidat chapter
		% 	introduced -> described (já asi těžko něco představil že ...)
		% 	hloubka obsahu?
		% 	opravdu dosmažit všechny 4 trace acq tooly
		% 	kam až číslovat? \subsubsection? změní \setcounter{secnumdepth}{4}
		% 	note vs. remark
		% 	XOR?
		% 	referencování na SCA algoritmy, neni lepší je pojmenovat?
		% 	
		
		
		%%% PŘEDMLUVA
		\cleardoublepage\phantomsection   % cleardoublepage a phantomsection udělají že odkaz z obsahu vede NAD nadpis, jinak vede POD
\addcontentsline{toc}{chapter}{Preface}
\chapter*{Preface}

% PŘEDMLUVA !!! se tomu řiká u nás ...
%~ význam WBC, kde se to vzalo, co za útok že tam chci použít, musí to bejt poutavý!
% o čem ta práce vlastně bude ve smyslu uvedení do problému, jak/proč ta práce vznikla?

%~ Motivation for whitebox crypto, whitebox crypto history, milestones   % Bos et al. to hezky shrnul
	%~ from black-box ... gray-box ... white-box
	%~ malicious software on your device
	%~ cryptographic keys of third-parties
	%~ use cases: DRM
	%~ modeled as if adversary had full control over execution environment of the device running the implementation
	%~ everything allowed: perform static analysis, read/alter memory (fault injection)
	%~ Chow et al. \cite{chow2003white}
	%~ program obfuscation impossibility results
	%~ goal of wbc = embed the secret key to the implementation s.t. it is difficult to extract it even when the source code is available
	%~ note the difference from obfuscation, in practice could be utilized as well
	%~ conjecture by Chow: wb implementations provide no long-term defense against attacks
	%~ symmetric crypto only, no published results on pubkey crypto
	%~ all published approaches have been broken
	%~ software vendors then disobey Kerckhoff's principle (cite) i.e. use secret design
	%~ Bos et al. exploit leaked information instead, no need for exact knowledge of design nor reverse engineering
	%~ perfect wb implementation should resist all existing and future side-channel attacks (C. Delerablee, T. Lepoint, P. Paillier, and M. Rivain. White-box security notions for symmetric encryption schemes. In T. Lange, K. Lauter, and P. Lisonek, editors, SAC 2013, volume 8282 of LNCS, pages 247–264, Burnaby, BC, Canada, Aug. 14–16, 2014. Springer, Berlin, Germany.)
	%~ ... Bos et al. utilize software traces for this purpose
	%~ The goal of this thesis is to reproduce the attack and make it easy for future researchers and analysis.

\section*{Work overview}
	
	\begin{description}
		\item[Chapter \ref{chap:intro}.] In this chapter we describe basic whitebox crypto techniques\ldots
		\item[Chapter \ref{chap:results}.] This chapter presents our practical results together with\ldots
	\end{description}

% ====================================================================
% ====================================================================
% ====================================================================

A table for fun follows:
\begin{center}
\begin{tabular}{|| c || c | c | c ||}
	\hline\hline
	~ & \multicolumn{2}{c|}{\bf $n\times n$ squares} & {\bf Computational} \\
	~ & \multicolumn{1}{c}{Lower bound} & \multicolumn{1}{c|}{Upper bound} & {\bf Power}\\
	\hline
	2D & \multicolumn{2}{c|}{See see} & Universal \\
	\hline
	3D & see & see & Universal \\
	\hline
	2D & see & see & Unknown \\
	\hline\hline
\end{tabular}
\end{center}

Just for fun:
\begin{figure}[H]
\begin{center}
	\includegraphics[width=0.4\textwidth]{./figures/fig.pdf} % {šířka v mm}/370 je to původně
	\caption{Remove me, please.}
	\label{fig:3color}
\end{center}
\end{figure}

Or a multifigure:
\begin{figure}[H]
\begin{center}
	\begin{subfigure}[b]{0.45\textwidth}
		\includegraphics[width=\textwidth]{./figures/fig.pdf}
		\caption{One.}
		\label{fig:one}
	\end{subfigure}
	\begin{subfigure}[b]{0.45\textwidth}
		\includegraphics[width=\textwidth]{./figures/fig.pdf}
		\caption{Two.}
		\label{fig:two}
	\end{subfigure}
	\caption{One and Two.}
	\label{fig:onetwo}
\end{center}
\end{figure}

		
		%%% KAPITOLA 1
		\chapter{Introduction to Whitebox Cryptography}
\label{chap:intro}

% Goal of crypto?
% This chapter will be organized as follows: ...

\section{Introduction to Cryptography}
\label{sec:introcrypto}

In cryptography, the most typical situation is where Alice wants to communicate with Bob confidentially over an insecure channel, while Eve may eavesdrop on the channel. In the beginning, let us suppose that Alice and Bob share a certain portion of secret information referred to as {\em secret key}, and both control a secure execution environment, where they intend to run their ciphering algorithm, so that Eve can only observe what they send over the channel.

\begin{note}
\label{note:threat}
	A list of such attacker's abilities is referred to as {\em threat model}. Note that in the previous threat model, the attacker is quite weak -- she can only read the traffic. If she were for instance able to amend the traffic, Alice and Bob would need to introduce other features to their protocol, so that they could detect forged messages.
	
	Before we begin with introducing any means of security, we always need to identify the threat model first. For now, we will only consider eavesdropping.
\end{note}

Our first intuitive requirements on Alice's and Bob's encryption algorithm could be expressed as follows: no matter what plaintext Alice and Bob encrypt and send over the channel, Eve shall {\em not} be able to
\begin{enumerate}
	\item recover any portion of any plaintext, \label{item:plainrecov}
	\item recover the secret key. \label{item:keyrecov}
\end{enumerate}
The item~\ref{item:plainrecov} requirement is also referred to as {\em confidentiality}. Note that the item~\ref{item:keyrecov} requirement -- key recovery -- would lead to plaintext recovery as well, on the other hand, not necessarily the other way around. Let us give an example of a primitive cipher.

\begin{example}
\label{ex:shift}
	{\em Shift cipher} \cite{menezes1996handbook} is an ancient cipher already used by Julius Caesar and it works as follows: it inputs a string of letters from a finite alphabet $\cal A$, maps its letters bijectively to the numbers from $0$ to $|{\cal A}|-1$ and adds a constant $c\in\{\atob{0}{|{\cal A}|-1}\}$ modulo $|\cal A|$ to each number. Finally, it maps the numbers back to letters yielding a string ``shifted'' by $c$, where $c$ can be interpreted as an element of $\cal A$ as well. Decryption can be achieved simply by substracting $c$ from the ciphertext in a similar way.
	
	Shift cipher can be easily broken by pen and paper only, since there are as many keys as letters in the alphabet, typically $26$ English letters. We simply try them all until the plaintext makes sense. Such approach is referred to as {\em brute force attack} or {\em exhaustive search}.
\end{example}
\nomenclature[01]{$\lvert S \rvert$}{Cardinality of the set $S$}

As we could have seen in the previous example, shift cipher is not very secure cipher. Therefore we would like to formalize the concept of a cipher and its security.


% ==============================================================================
% ===   S Y M M E T R I C   A N D   A S Y M M E T R I C   C I P H E R        ===
% ==============================================================================

\subsection{Symmetric And Asymmetric Cipher}

There are two large families of ciphers -- symmetric and asymmetric ones. Shift cipher, introduced in Example~\ref{ex:shift}, belongs to the symmetric family, which follows from the fact that it uses the same key for both encryption and decryption, while asymmetric ciphers use distinct keys for encryption and decryption. Asymmetric ciphers are also referred to as public key ciphers, but we will only consider symmetric ciphers in this thesis; a definition follows.

\begin{defn}[Symmetric Cipher]
\label{def:symcipher}
	Let $K$ denote the key space, $M$ the message space and $C$ the ciphertext space. {\em Symmetric cipher} is a pair of efficiently computable (possibly randomized) functions $(E,D)$, $E:K\times M\rarr C$, $D:K\times C\rarr M$, such that $\forall m\in M$ and $\forall k\in K$ it holds   % for fun: \Pr\left(D(E(l,k),k)=l\right) = 1 .
	\begin{equation}
	\label{eq:correct}
		\Pr\Bigl(D\bigl(k,E(k,m)\bigr)=m\Bigr) = 1 .
	\end{equation}
\end{defn}
\nomenclature[02]{$\Pr(\omega)$}{Probability of the event $\omega$}

\begin{note}
	Going back to Example~\ref{ex:shift}, the key space is $\{\atob{0}{|{\cal A}|-1}\}\sim\cal A$, the message and ciphertext spaces are ${\cal A}^*$ which stands for the set of all strings over the alphabet $\cal A$.
\end{note}

According to the previous definition, we do not seem to have a definition of a proper cipher at all -- the only requirement is corectness (cf.\ Equation~\ref{eq:correct}), but it totally lacks any security requirement. As we will see later, it is actually quite complicated to define ``security'' properly.


% ==============================================================================
% ===   I N F O R M A T I O N - T H E O R E T I C   A P P R O A C H          ===
% ==============================================================================

\subsection{Information-Theoretic Approach}

Shannon's groundbreaking work \cite{shannon1949mathematical} on the mathematics of communication gives us a tool -- information-theoretic approach. We can then define a secure cipher as a cipher, for which it holds that its ciphertexts carry no information about respective plaintexts. Note that this can be rewritten in a more friendly form, a definition follows.

\begin{defn}[Perfectly Secure Cipher]
\label{def:perfsec}
	Let $(E,D)$ be a symmetric cipher. Then it is called {\em perfectly secure} if $\forall m_0,m_1\in M$ such that $|m_0| = |m_1|$ and $\forall c\in C$ it holds $\Pr\bigl(E(k,m_0)=c\bigr) = \Pr\bigl(E(k,m_1)=c\bigr)$, where $k$ is uniformly random in $K$, denoted by $k\unirand K$.
\end{defn}
\nomenclature[03]{$k\unirand K$}{$k$ is taken uniformly random from the set $K$}

\begin{note}
\label{note:indist}
	In other words, the attacker is not able to {\em distinguish} encryption of $m_0$ from encryption of $m_1$ by {\em any} means.
\end{note}

There is a cipher that is perfectly secure, see the following example.

\begin{example}[Vernam Cipher]
	Let $K = M = C = \{0,1\}^n$ and $k\unirand K$. Given a plaintext $m$, {\em Vernam cipher} applies bitwise XOR of the key $k$ and the plaintext $m$ yielding a ciphertext $c$, i.e., $c = m\xor k$.
\end{example}
\nomenclature[04]{$a\xor b$}{Bitwise XOR of two bitstrings $a$ and $b$}

\begin{note}
	Vernam cipher is also referred to as {\em One Time Pad} (OTP). Note that it is crucial that each key is only used once, otherwise we could attack the cipher easily. Indeed, given two ciphertexts using the same key, i.e., $c_0=m_0\xor k$ and $c_1=m_1\xor k$, one can compute $c_0\xor c_1 = m_0\xor m_1$. This can be practically attacked using some apriori knowledge about the plaintexts (e.g., encoding -- ASCII, format and language of the plaintext -- XML and English, and so on), which reduces the amount of possibilities and typically leads to a complete plaintext recovery.
\end{note}

Perfect security seems to be exactly what we want from a secure cipher, but the following theorem shows that it implies certain inconveniences.

\begin{thm}
\label{thm:kgeqm}
	Let $(E,D)$ be a perfectly secure cipher. Then $|K| \geq |M|$.
\end{thm}

\begin{note}
\label{note:intent}
	This consequence of perfect security apparently goes against our intention -- we would like to establish a shared secret key of a limited length and then, using the key, be able to encrypt as much information as we want. For this reason, we need to weaken our security assumption. Note that previously, there was no assumption on attacker's computing power -- the requirement was actually to hide the  plaintext from {\em any} attacker, i.e., including an attacker with unlimited computing power.
\end{note}


% ==============================================================================
% ===   C O M P L E X I T Y - T H E O R E T I C   A P P R O A C H            ===
% ==============================================================================

\subsection{Complexity-Theoretic Approach}

We would like to utilize the fact that the attacker typically does not possess unlimited computing power, but she can only solve probabilistic polynomial problems.

\begin{note}
\label{note:polyinterms}
	There arises a natural question: {\em Probabilistic polynomial in terms of what?} In cryptography, there is a meter for that referred to as {\em security parameter}, which is typically the key-length. In the following text and for this purpose, the expression {\em effective} will be used frequently.
\end{note}

With this assumption, we can weaken our requirement on ciphertext from ``carries no information'' to ``is computationally impossible to gain some information''. In a manner analogous to Note~\ref{note:indist}, we would like to make impossible to distinguish which plaintext has been encrypted, now considering an effective attacker only. Let us present this idea in the following game.

\begin{game}
\label{game:semsec}
	There are two parties -- a {\em challenger} and an {\em adversary}. First, the challenger chooses a random key $k\unirand K$ and a random bit $b\unirand\{0,1\}$. Then the adversary sends two {\em equally long} plaintexts $m_0$, $m_1$ of her choice to the challenger. The challenger encrypts $m_b$ (according to $b$ she has chosen) and sends the resulting ciphertext $c$ back to the adversary. The adversary is allowed to perform several such queries while the challenger keeps constant key $k$ and bit $b$; the adversary is only limited by polynomial time.
	
	The aim of the adversary is to distinguish effectively which $b$ has been chosen by the challenger; see Figure~\ref{fig:semsecgame} depicting this game.
\end{game}

\begin{note}
\label{note:semsec}
	Resistance of a cipher to this kind of attack is referred to as {\em semantic security}.
\end{note}

\begin{figure}[H]
\begin{center}
	\includegraphics{./figures/game_semsec/game_semsec-1.pdf}
	\caption{Game~\ref{game:semsec} depicted.}
	\label{fig:semsecgame}
\end{center}
\end{figure}

\begin{note}   %_% NEW
\label{note:ambigencr}
	The adversary is intentially allowed to provide the same plaintext over and over again. However, it follows that the challenger has to encrypt the same plaintext differently each time, otherwise there would exist a trivially winning adversary.
	
	Indeed, such an adversary would pick three distinct plaintexts $m_0$, $m_1$, $m_2$ and send $(m_0,m_1)$ to the challenger, while she would get back some ciphertext $c_a$. In her second query, the adversary would send $(m_0,m_2)$ to the challenger and receive another ciphertext, say $c_b$. If encryptions of $m_0$ were equal, she could decide with certainty:
	\begin{itemize}
		\item $b' = 0$ if $c_a = c_b$, or
		\item $b' = 1$ otherwise,
	\end{itemize}
	and win the game.
\end{note}

\begin{note}   %_% NEW
\label{note:randomize}
	In order to achieve semantic security of a cipher, the cipher {\em must} encrypt the same plaintext differently each time. This can be achieved using randomization or counters, both approaches are practically used.
\end{note}


\subsubsection{Definition of Semantic Security}

Let us introduce some notion first.

\begin{defn}[Negligible Function]
\label{def:neglfunc}
	Let $\epsilon:\N\rarr\R$. $\epsilon$ is called {\em negligible} if $\forall d\in\N$ $\exists \lambda_0\in\N$ such that $\forall \lambda>\lambda_0$ it holds $\epsilon(\lambda)\leq\frac{1}{\lambda^d}$. In the opposite case, $\epsilon$ is called {\em non-negligible}.
\end{defn}
\nomenclature[06]{$\N$}{Positive integers}
\nomenclature[07]{$\R$}{Real numbers}

\begin{defn}[Overwhelming Function]
\label{def:overwh}
	Let $\epsilon:\N\rarr\R$. $\epsilon$ is called {\em overwhelming} if $1-\epsilon$ is negligible.
\end{defn}

\begin{note}
\label{note:neglconst}
	In practice, {\em negligible} is commonly used with concrete constants as well: $\epsilon<\frac{1}{2^{80}}$ is considered negligible, i.e., events with such probability are not considered to occur during our lives. On the other hand, $\epsilon>\frac{1}{2^{30}}$ is considered non-negligible, i.e., it is likely to observe an event with such probability.
\end{note}

\begin{defn}[Oracle]
\label{def:oracle}
	Let $F:X\rarr Y$. {\em Oracle evaluating $F$} is a device that, given $x\in X$, evaluates and returns $F(x)$ in unit time; denoted by ${\cal O}_F$.
\end{defn}

\begin{defn}[Oracle with Limited Access]   %_% NEW
\label{def:saoracle}
	Let $F:X\rarr Y$ and $n\in\N$. {\em Oracle with limited access evaluating $F$} is an oracle evaluating $F$, where the number of queries is limited to $n$; denoted by ${\cal O}^{(n)}_F$.
\end{defn}

\begin{defn}[Statistical Test]
\label{def:stattest}
	Let $F:X\rarr Y$. {\em Statistical test} is a (possibly randomized) effective algorithm $A$, which only has an access to an oracle evaluating $F$, and which outputs $0$ or $1$; denoted by $A({\cal O}_F)$.
\end{defn}

\begin{note}
	Statistical test will be also referred to as {\em adversary}.
\end{note}

\begin{defn}[Advantage]
\label{def:advant}
	Let $\cal F$, $\cal G$ be two distinct subsets of the set of all functions from $X$ to $Y$, and let $A$ be an adversary. We define {\em advantage} of adversary $A$ as
	\[
		\Adv(A,{\cal F},{\cal G}) = \left| \Pr\limits_{F\unirand {\cal F}}\Bigl(A({\cal O}_F)=1\Bigr) - \Pr\limits_{G\unirand {\cal G}}\Bigl(A({\cal O}_G)=1\Bigr) \right| .
	\]
\end{defn}

\begin{defn}[Advantage with Limited Access]   %_% NEW
\label{def:saadvant}
	Let $n\in\N$. {\em Advantage with limited access} is defined as advantage, where oracles with access limited to $n$ are used instead; denoted by ${\Adv}^{(n)}$.
\end{defn}

In other words, advantage tells us how likely adversary $A$ is able to distinguish a random function from one set of functions from a random function from another set, based on its input/output behavior only. According to these sets, advantage can be ``customized'' for a specific purpose. Let us define such advantage for an adversary tempting to distinguish which of her plaintexts has been encrypted, i.e., tells us how likely she is to win the Game~\ref{game:semsec}.

\begin{defn}[Semantic Security Advantage]
\label{def:ssadvant}
	Let $E: K\times M\rarr C$ be an encryption function of a symmetric cipher, ${\cal E}_b = \{E_b:M^2\rarr C \mid E_b(m_0,m_1) = E(k,m_b), k\in K\}$ for $b=0,1$, and $A$ an adversary. We define {\em semantic security advantage} of adversary $A$ as
	\[
		\AdvSS(A,E) = \Adv(A,{\cal E}_0,{\cal E}_1) .
	\]
\end{defn}

Semantic security advantage gives us a reasonable meter for formalization of Game~\ref{game:semsec}. Let us use it to define another notion of security, cf.\ Definition~\ref{def:perfsec} (Perfectly Secure Cipher).

\begin{defn}[Semantically Secure Cipher]
\label{def:semsec}
	Let $(E,D)$ be a symmetric cipher. Then it is called {\em semantically secure} if, for {\em any} adversary $A$, the semantic security advantage $\AdvSS(A,E)$ is negligible.
\end{defn}

Semantically secure cipher therefore cannot be broken by any probabilistic-polynomial adversary with overwhelming probability. On the other hand, it might be broken by an adversary with unlimited computing power -- unlike a perfectly secure cipher that cannot be broken by any means; see Note~\ref{note:indist}.

The benefit of semantic security over perfect security is that semantic security does not imply any inconvenient consequence analogous to Theorem~\ref{thm:kgeqm}, i.e., addresses our intention as stated in Note~\ref{note:intent}, and practically provides similar security.

Let us now give a nice real-world example, how things can go wrong if semantic security is not complied, even though the design follows common sense.

\begin{example}
\label{ex:crime}
	In 2012, a combination of settings in TLS\footnote{Transport Layer Security, a cryptographic protocol \cite{dierks2008transport}.} allowed researchers to hijack a web session by cracking encrypted cookies \cite{goodin2012crime}. The problem arised when compression was followed by encryption, which actually makes sense. Let us have a look from theoretical point of view.
	
	Remind the condition of Game~\ref{game:semsec} -- it requires two plaintexts of equal length. With this assumption, the respective ciphertexts were indistinguishable. However, if the plaintext is compressed first, it can result in differently-sized inputs to the subsequent (semantically secure) cipher! This obviously harms the assumption of inputs of equal length, and may allow the adversary to distinguish the ciphertexts. Note that only very little queries (around six) were needed to attack one byte of an encrypted cookie in the practical attack!
	
	This vulnerability was given a CVE\footnote{Common Vulnerabilities and Exposures, a standard for information security vulnerability names, see \link{cve.mitre.org}.} code CVE-2012-4929 and a human-readable name ``CRIME'' (Compression Ratio Info-leak Made Easy).
\end{example}

\subsubsection{Shannon's Principles}
	
	Before we present how semantic security is practically achieved, let us introduce a principle that is widely followed in cipher design.
	
	\begin{princ}[Shannon]
	\label{pri:shannon}
		As outlined in another Shannon's work \cite{shannon1949communication}, secure ciphers are suggested to employ two basic principles -- {\em confusion} and {\em diffusion}.
		\begin{description}
			\item[Confusion.] The relation between information in plaintext or key and information in respective ciphertext shall be very complex, e.g., should not belong to a specific subset of mappings, especially to linear ones.
			\item[Diffusion.] Small piece of information in plaintext or key shall be spread in a long piece of respective ciphertext, e.g., one bit of plaintext or key should change a large portion of respective ciphertext.
		\end{description}
	\end{princ}
	
	Semantic security is usually achieved using smaller building blocks, which are required to have certain properties -- they especially make use of confusion and diffusion. The following section describes a building block of one of possible approaches.


% ==============================================================================
% ===   B L O C K   C I P H E R                                              ===
% ==============================================================================

\subsection{Block Cipher}

There are two main approaches how to construct a candidate semantically secure cipher. They make use of
\begin{enumerate}
	\item a {\em stream cipher}, or
	\item a {\em block cipher},
\end{enumerate}
respectively. Stream ciphers achieve semantic security by proper initialization, on the other hand, block ciphers use initialization together with {\em modes of operation}, e.g., {\em Cipher Block Chaining} (CBC) or {\em Counter Mode} (CTR). Modes of operation will not be covered in this thesis; only block ciphers and their desired properties will be described, a bunch of definitions follows.

\begin{defn}[Pseudorandom Permutation]
\label{def:prp}
	Let $E: K\times X\rarr Y$. We call $E$ a {\em pseudorandom permutation} (PRP) if there exists an efficient algorithm that evaluates $E$, the function $E(k,\cdot): X\rarr Y$ is a bijection and there exists an efficient algorithm to compute $E^{-1}(k,\cdot)$, for every $k\in K$.
\end{defn}

\begin{note}
	Pseudorandom permutation is usually referred to as {\em block cipher}. Note that the previous definition is in some sense very similar to Definition~\ref{def:symcipher} -- it does not introduce any security requirements either.
\end{note}

Pseudorandom permutations will play the main role from now and $k$ will play the role of their key. In order to achieve semantic security of certain constructions using a PRP, we need the PRP to behave indistinguishable from a truly random bijection. Let us demonstrate a possible approach on a game similar to Game~\ref{game:semsec}.

\begin{game}
\label{game:prp}
	There are two parties -- a {\em challenger} and an {\em adversary}. Let $E:K\times X\rarr Y$ be a PRP. First, the challenger chooses a random key $k\unirand K$, a random bijection $F:X\rarr Y$ and a random bit $b\unirand\{0,1\}$. Then the adversary sends $x\in X$ of her choice to the challenger. The challenger applies $E(k,\cdot)$ or $F$ if $b=0$ or $1$, respectively, and sends the result back to the adversary. The adversary is allowed to perform several such queries.
	
	The aim of the adversary is to distinguish effectively which $b$ has been chosen by the challenger.
\end{game}

Before we give a formal definition of a secure PRP, we define adversary's advantage in the previous game in a manner analogous to Definition~\ref{def:ssadvant}.

\begin{defn}[PRP Advantage]
\label{def:prpadvant}
	Let $E: K\times X\rarr Y$ be a PRP, $\cal B$ the set of all bijections from $X$ to $Y$, and $A$ an adversary. We define {\em PRP advantage} of adversary $A$ as
	\[
		\AdvPRP(A,E) = \Adv\Bigl(A,\bigl\{E(k,\cdot) \bigm| k\in K\bigr\},{\cal B}\Bigr) .
	\]
\end{defn}

PRP advantage tells us how likely adversary $A$ is able to distinguish a PRP with a random key from a truly random bijection based on its input/output behavior only. The definition of secure PRP follows.

\begin{defn}[Secure PRP]
\label{def:secprp}
	Let $E: K\times X\rarr Y$ be a PRP. Then it is called {\em secure PRP} if, for {\em any} adversary $A$, the PRP advantage $\AdvPRP(A,E)$ is negligible.
\end{defn}

In other words, given a random key, PRP is secure if it is computationaly indistinguishable from a truly random bijection. 

An example covering most introduced concepts follows.

\begin{example}
	Let $E:\{0,1\}^{128}\times\{0,1\}^{128}\rarr\{0,1\}^{128}$ be a PRP defined as
	\[
		E(k,x) = \bin\bigl((k)_2 + (x)_2\mod{2^{128}}\bigr)
	\]
	where $\bin(\cdot)$ stands for number's $128$-bit binary representation and $(\cdot)_2$ stands for a number represented by given binary string.
	
	Now let us construct an adversary $A$ as follows: it generates two arbitrary but distinct $x_0,x_1\in\{0,1\}^{128}$, and feeds them to an oracle computing either PRP $E$, or a truly random bijection. The oracle returns $y_0$ and $y_1$, respectively, and the adversary compares values $(y_i)_2 - (x_i)_2 \mod{2^{128}}$ for $i=0,1$ with each other. If these values are equal, $A$ returns $1$, otherwise it returns $0$.
	
	Note that if the oracle computes PRP $E$, then these values are always equal. Otherwise -- i.e., the oracle computes a truly random bijection -- the values are distinct with overwhelming probability. Therefore the PRP advantage of such adversary is
	\[
		\AdvPRP(A,E) = \left| \Pr\limits_{k\unirand \{0,1\}^{128}}\Bigl(A\bigl({\cal O}_{E(k,\cdot)}\bigr)=1\Bigr) - \Pr\limits_{F\unirand {\cal B}}\Bigl(A({\cal O}_F)=1\Bigr) \right| = 1 - \epsilon ,
	\]
	where $\cal B$ stands for the set of all bijections on $\{0,1\}^{128}$ and $\epsilon$ is negligible.
	
	It follows that such PRP is totally insecure, since there exists an adversary with overwhelming advantage, but only a negligible advantage is allowed.
\end{example}

\begin{note}   %_% NEW
\label{note:secbyobsc}
	In the previous example, the adversary exploits the knowledge of the block cipher's internal structure -- she knows everything but the key. One could have an idea to hide the structure from the adversary, but this is actually a {\em very} bad idea. Once our construction leaks, we need to deploy a new ciphering algorithm confidentially, but we have no way to do so. Moreover, note that a secure PRP was required to be resistant to {\em all} adversaries, hence hiding your design has absolutely no theoretical support.
	
	We shall rather use a secure publicly known block cipher and only rely on the key. Once our key is compromised, we simply establish a new key. Also note that our key is only to be protected within our device, but the cryptographic algorithm would have to be protected globally.
	
	Such approach -- hiding the design of a cipher -- is referred to as {\em security by obscurity} and shall be strictly avoided.
\end{note}

This idea has already been addressed in 1883 by Kerckhoffs \cite{auguste1883cryptographie}. It is usually referred to as {\em Kerckhoffs' Principle} and its six rules follow (originally French, English translation taken from \cite{petitcolas2016kerckhoffs}):
\begin{princ}[Kerckhoffs]
	~
	\begin{enumerate}
		\item The system must be substantially, if not mathematically, undecipherable;
		\item The system must not require secrecy and can be stolen by the enemy without causing trouble;\label{item:kerck}
		\item It must be easy to communicate and retain the key without the aid of written notes, it must also be easy to change or modify the key at the discretion of the correspondents;
		\item The system ought to be compatible with telegraph communication;
		\item The system must be portable, and its use must not require more than one person;
		\item Finally, given the circumstances in which such system is applied, it must be easy to use and must neither stress the mind or require the knowledge of a long series of rules.
	\end{enumerate}
\end{princ}
Hence harming rule~\ref{item:kerck} from the previous list leads to security by obscurity as introduced in Note~\ref{note:secbyobsc}.


% ==============================================================================
% ===   O U T P U T   I N D I S T .   O F   A   S E C U R E   P R P          ===
% ==============================================================================

\subsection{Output Indistinguishability of a Secure PRP}

Now we know what a secure PRP is (or a secure block cipher) -- it is based on Game~\ref{game:prp}, where the aim of the adversary is to distinguish a PRP from a truly random bijection. The obvious question is, what happens if we put a secure PRP into Game~\ref{game:semsec}, where the adversary tempts to distinguish which of her two inputs has been processed.

Note that PRP {\em is} a cipher (cf.\ Definition~\ref{def:symcipher} and \ref{def:prp}), so it makes sense. Indeed, given a PRP, there exists an effective algorithm evaluating it and its inverse, and also Equation~\ref{eq:correct} holds.

\begin{note}   %_% NEW
\label{note:singleaccess}
	PRP has no ambiguity, hence we will need to either limit the amount of adversary-challenger queries to one, or, equivalently, change the key after each query. If we did not do so, then there would obviously exist a trivially winning adversary as outlined in Note~\ref{note:ambigencr}.
\end{note}

Let us modify Game~\ref{game:semsec} in a manner described in Note~\ref{note:singleaccess} for the purpose of testing PRPs.

\begin{game}
\label{game:semsecprp}
	There are two parties -- a {\em challenger} and an {\em adversary}. Let $E:K\times X\rarr Y$ be a PRP. First, the challenger chooses a random key $k\unirand K$ and a random bit $b\unirand\{0,1\}$. Then the adversary sends two equally long plaintexts $m_0$, $m_1$ of her choice to the challenger. The challenger computes $c = E(k,m_b)$ and sends it back to the adversary. The adversary is limited to one query only.
	
	The aim of the adversary is to distinguish effectively which $b$ has been chosen by the challenger.
\end{game}

Based on the previous game, let us define {\em PRP semantic security advantage} and {\em semantically secure PRP} (cf.\ Game~\ref{game:semsec}, Definition~\ref{def:ssadvant} and \ref{def:semsec}). Note that the difference is basically limited oracle access.

%_% NEW
% BEGIN
\begin{defn}[PRP Semantic Security Advantage]
\label{def:prpssadvant}
	Let $E: K\times X\rarr Y$ be a PRP, ${\cal E}_b = \{E_b:X^2\rarr Y \mid E_b(x_0,x_1) = E(k,x_b), k\in K\}$ for $b=0,1$, and $A$ an adversary. We define {\em PRP semantic security advantage} of adversary $A$ as
	\[
		\AdvPRPSS(A,E) = {\Adv}^{(1)}(A,{\cal E}_0,{\cal E}_1) .
	\]
\end{defn}

\begin{defn}[Semantically Secure PRP]
\label{def:semsecprp}
	Let $E: K\times X\rarr Y$ be a PRP. Then it is called {\em semantically secure PRP} if, for {\em any} adversary $A$, the PRP semantic security advantage $\AdvPRPSS(A,E)$ is negligible.
\end{defn}
% END

The following theorem provides a connection of Game~\ref{game:prp} and \ref{game:semsecprp}. It claims that, given a PRP $E$ winning Game~\ref{game:prp}, i.e., $E$ is indistinguishable from a truly random bijection, $E$ wins Game~\ref{game:semsecprp} as well, i.e., no adversary can distinguish encryption of $m_0$ from encryption of $m_1$.

\begin{thm}
\label{thm:semsecprp}
	Let $E: K\times X\rarr Y$ be a secure PRP. Then it is also a semantically secure PRP.
\end{thm}
\begin{proof}
	We show that for every PRP semantic security adversary $A$, there exist PRP security adversaries $B_0$, $B_1$ such that
	\begin{equation}
	\label{eq:aleqbb}
		\AdvPRPSS(A,E) \leq \AdvPRP(B_0,E) + \AdvPRP(B_1,E) .
	\end{equation}
	The claim then easily follows -- if $E$ is a secure PRP, i.e., $\AdvPRP(B_0,E) + \AdvPRP(B_1,E)$ is negligible, then also $\AdvPRPSS(A,E)$ is negligible. Hence $E$ is a semantically secure PRP.
	
	Let us first denote $W_b$ the event when adversary $A$ answers $1$ while the true value of the PRP semantic security game is $b$, for $b=0,1$. In this notation, it holds
	\begin{equation}
	\label{eq:prpssae}
		\AdvPRPSS(A,E) = \bigl| \Pr(W_1) - \Pr(W_0) \bigr| .
	\end{equation}
	Let us further denote $R_b$ the same event with only difference -- the PRP semantic security game uses a truly random bijection instead. It holds
	\begin{equation}
	\label{eq:rr0}
		\bigl| \Pr(R_1) - \Pr(R_0) \bigr| = 0 ,
	\end{equation}
	since $b\unirand\{0,1\}$ and both $R_b$ provide a truly random bijection independent from $b$.
	
	Now we construct adversaries $B_0$, $B_1$ distinguishing PRPs from truly random ones while we make use of adversary $A$. Let us construct $B_0$ first. We ask $A$ to provide $m_0$ and $m_1$, then we feed the PRP oracle with $m_0$ obtaining $c$ as a result. Note that either $c=E(k,m_0)$ or $c=F(m_0)$ according to $b$, where $F$ is a truly random bijection. We provide this $c$ to adversary $A$, and output what $A$ outputs. The advantage of adversary $B_0$ is
	\begin{equation}
	\label{eq:prpb0}
		\AdvPRP(B_0,E) = \left| \Pr(W_0) - \Pr(R_0) \right| .
	\end{equation}
	We construct $B_1$ in a similar manner (i.e., choosing $m_1$ instead) and get
	\begin{equation}
	\label{eq:prpb1}
		\AdvPRP(B_1,E) = \left| \Pr(W_1) - \Pr(R_1) \right| .
	\end{equation}
	Combining Equation~\ref{eq:prpssae}, \ref{eq:rr0}, \ref{eq:prpb0} and \ref{eq:prpb1}, we get
	\begin{align*}
		\AdvPRPSS(A,E) &= \left| \Pr(W_1) - \Pr(W_0) \right| \leq \\
		&\leq \left| \Pr(W_1) - \Pr(R_1) \right| + \left| \Pr(W_0) - \Pr(R_0) \right| = \\
		&= \AdvPRP(B_1,E) + \AdvPRP(B_0,E) ,
	\end{align*}
	which finally proves Equation~\ref{eq:aleqbb}, hence the theorem, too
\end{proof}

\begin{note}
	The previous theorem does not work the other way around, i.e., a semantically secure PRP does not need to be a secure PRP. Indeed, let $E: K\times X\rarr Y$ be a secure PRP, and let us define $E'(k,x) = E(k,x)\|1$, where $\|$ stands for string concatenation. $E'$ is obviously semantically secure (it remains hard to distinguish encryptions of dinstinct plaintexts), but it can be easily distinguished from a truly random bijection, since it always ends with $1$. The advantage of such adversary would be $\AdvPRP(A,E') = \nicefrac{1}{2}$ which is clearly non-negligible.
\end{note}
\nomenclature[05]{$a \parallel b$}{Concatenation of two strings $a$ and $b$}

The previous theorem is especially important for the idea behind its proof, which is used in many other proofs related to provable-security. In this proof, we constructed adversaries $B_0$, $B_1$, which internally employed adversary $A$, and then we used advantage of $A$ to estimate advantages of $B_0$ and $B_1$.


% ==============================================================================
% ===   P R O V A B L E   V S .   R E A L - W O R L D   S E C U R I T Y      ===
% ==============================================================================

\subsection{Provable vs. Real-World Security}
\label{sec:provable}

Provable security is one aspect, the other is practical construction of a cipher. Note that there is no provably secure block cipher known.

If there were, it would actually imply that $\Pclass \subsetneq \NPclass$, which is an open millenium problem. Indeed, if $\Pclass = \NPclass$ and $E$ were a PRP, then the adversary could, given certain (polynomial) amount of plaintext-ciphertext pairs $(m_i,c_i)$, non-deterministically find a key $k$, such that $c_i = E(k,m_i)$ for all $i$. Then she would simply answer ``truly random'' if no such $k$ were found, otherwise she says ``pseudorandom''. Here, the ``truly random'' answer is answered with certainty, while the ``pseudorandom'' answer is just a conjecture, but it has sufficiently large probability.
\nomenclature[08]{$\Pclass$}{Deterministic polynomial complexity class, i.e., the set of decision problems which can be solved by a deterministic Turing machine in polynomial time}
\nomenclature[09]{$\NPclass$}{Non-deterministic polynomial complexity class, i.e., the set of decision problems which can be solved by a non-deterministic Turing machine in polynomial time}

Instead, real-world block ciphers are {\em believed} to be secure \cite{menezes1996handbook} based on a simple fact: no known attack is far better than exhaustive search. One such block cipher is called {\em Advanced Encryption Standard} (AES) and will be of top interest in this thesis. AES will be described in detail in Section~\ref{sec:aes}.


% ==============================================================================
% ===   F R O M   B L A C K - B O X   T O   W H I T E - B O X                ===
% ==============================================================================

\subsection{From Black-Box to White-Box Attack Context}
\label{sec:bbtowb}

The other central topic of this thesis is {\em white-box attack context}. In order to explain what it is, let us begin with {\em black-box attack context}.

\subsubsection{Black-Box}

First of all, as stated in Note~\ref{note:threat}, we need to address attacker's abilities -- the threat model. In the black-box model, the attacker is granted an access to an oracle evaluating given encryption algorithm $E$ with a random secret key $k$, while the (primary) goal of the attacker is to recover the key. The attacker has absolutely no information but inputs and respective outputs, she cannot observe any internal routines of the encryption procedure, hence {\em black-box}.

\begin{note}
	There could be many other goals than just key recovery, e.g., guessing the ciphertext in advance, or decrypting certain/given ciphertext with non-negligible probability and so on. Actually whatever one can imagine that would be in principle possible with a PRP, but not with a truly random bijection.
\end{note}

Classical ciphers are designed under the assumption of black-box attack context, therefore their implementation must be extremely careful. The attacker could for example measure the time of encryption and exploit this information, which was not considered during the design of the cipher.

\subsubsection{Gray-Box}

Gray-box attack context is closer to reality, since we do not encrypt using oracles, but rather using some physical device, which runs a specific implementation of given cipher. In the gray-box model, cipher is thus supposed to be implemented and run, while the attacker can observe some information related to cipher's intermediate results as well.

Such information leaks can be modelled as a set of distributions conditioned by intermediate results and are referred to as {\em leakage model}. There is an area of theoretical cryptography studying leakages -- {\em Leakage Resilient Cryptography}. A practical attack exploiting such leaks will be described in Chapter~\ref{chap:attack} in Section~\ref{sec:side}.

\subsubsection{White-Box}

Finally, there is the most extreme case where the attacker has full control over the execution environment -- the {\em white-box attack context} (WBAC). Full control means that the attacker can step the algorithm and directly observe and alter intermediate results in the memory, or even skip, alter and insert instructions. Clearly, the key cannot be kept in its original form, it must be somehow hidden, but still usable for encryption.


% previously: Finally there is the most extreme case where the attacker has full control over the execution environment -- the white-box attack context. Full control means that the attacker can step the algorithm and directly observe and alter byproducts in the memory or even skip, alter and insert instructions. Clearly, the key cannot be kept in its original form, it must be somehow hidden but still usable for encryption.

\section{White-Box Cryptography}
	
	
	
%~ detailně popsat kde se vzala vč. citací, vývoj, hra na kočku a myš, impos. results
%~ že neni žádnej public key wb

% je to de facto threat model
%~ straightforward AES implementation is totally insecure if certain byproducts or even some information related to them is revealed. An attack exploiting such information leaks is covered in detail in Section \ref{sec:side}.

%!% přesun do přechozí sekce, přidat hru i před secure PRP
%~ Note that if a block cipher is secure then such advantage is also negligible.
% implementation can help, e.g. timing attacks
% RSA, an asymmetric cipher, could serve as a good example
% black box: like an oracle, everything holds. ciphers designed wrt black box, why wb? what it is? ...

%~ \subsubsection{Black-Box Model}
	
	%~ approach: feistel, iterated; diffusion vs. confusion


\section{Advanced Encryption Standard}
\label{sec:aes}

In 1997, Curtin and Dolske were the first to publicly announce (in \cite{curtin1998brute}) that they have cracked Data Encryption Standard (DES) \cite{fips1977des} -- a symmetric block cipher which had been in use for $20$ years. In the same year, the National Institute of Standards and Technology (NIST) started a selection process for a new encryption standard. During the process, DES was partially replaced with Triple-DES described in ANSI draft X9.52 \cite{americantripple}.

In 1998, NIST chose $15$ candidates and asked cryptology community for help with analyzing the candidates. After a careful analysis, NIST proposed Rijndael\footnote{From authors' names: Vincent Rijmen and Joan Daemen.} \cite{daemen1999aes} as the new standard called Advanced Encryption Standard (AES) \cite{fips2001aes} on November 26, 2001. The standard specifies the Rijndael algorithm with block-length of $128$ bits and key-lengths of $128$, $192$ and $256$ bits. Note that the original Rijndael algorithm \cite{daemen2013rijndael} describes more block- and key-length variants. On the other hand, we will only consider its $128$-bit variant in the following text, unless stated otherwise.

Today, 15 years later, AES is used worldwide and, although extensively analysed, is still beleived to be secure enough.


% ======================================================================
% ===   A E S   A L G O R I T H M   P R E R E Q U I S I T I E S      ===
% ======================================================================

\subsection{AES Algorithm Prerequisities}

AES encryption\footnote{Decryption will be treated separately at the end of this section.} consists of $10$ rounds, each round (except for the last one) consists of $4$ different operations which will be described later. All of these operations work with $128$-bit words, let us call them {\em states}.

Before we start, let us consider state differently: we split its $128$ bits into $16$ bytes denoted by $A_{0,0}, A_{1,0}, A_{2,0}, A_{3,0}, A_{0,1}, \ldots, A_{3,1}, \ldots, A_{3,3}$, respectively, obtaining a $4\times 4$ array where $A_{i,j}$ stands for an element in $i$-th row and $j$-th column.

Let us further denote $F = \GF(2)[x]/x^8+x^4+x^3+x+1$ which is a field of polynomials with integral coefficients reduced$\mod{2}$, these polynomials are then reduced$\mod(x^8+x^4+x^3+x+1)$. This field has $2^8$ elements and is usually referred to as {\em Rijndael's field} or simply $\GF(2^8)$.

\nomenclature{$\GF(p^n)$}{Galois Field with $p^n$ elements, $p$ prime}
\nomenclature{$F[x]$}{Ring of polynomials over $F$}

Each byte $A_{i,j} = a_7a_6\ldots a_0$, $a_i\in\GF(2)$, is then considered as an element of $F$ -- simply as $a_7x^7 + a_6x^6 + a_5x^5 + a_4x^4 + a_3x^3 + a_2x^2 + a_1x + a_0$.  Note that byte may also be written in hexadecimal notation and still considered as an element of $F$. See the following example with a few possible notations.
\begin{equation*}
\arraycolsep=.5em\def\arraystretch{1.5}
	01010001\,11010110\,\ldots\;\sim\;\texttt{51}\,\texttt{d6}\,\ldots\;\sim\;
	\begin{array}{|c|c|c|c|}
		\hline
		x^6 + x^4 + 1 & \ldots & \ldots & \ldots \\
		\hline
		x^7 + x^6 + x^4 + x^2 + x & \ldots & \ldots & \ldots \\
		\hline
		\ldots & \ldots & \ldots & \ldots \\
		\hline
		\ldots & \ldots & \ldots & \ldots \\
		\hline
	\end{array}
\end{equation*}
Definitions of the $4$ operations follow.

\subsubsection{SubBytes}
	
	The $\SubBytes$ operation is often referred to as {\em SBox} and denoted by $S$. It is defined byte-wise, i.e. $S: F \rightarrow F$. For $A\in F$, let $A'$ stand for
	\begin{itemize}
		\item $A^{-1}$, if $A \neq 0$,
		\item $0$, if $A = 0$.
	\end{itemize}
	Before we define $\SubBytes$ we must emphasize that the following multiplication is {\em not} performed in $F$ but in the ring $\GF(2)[x]$ instead. $\SubBytes$ is defined as follows:
	\begin{equation*}
		S(A) = (x^4 + x^3 + x^2 + x + 1)A' + x^6 + x^5 + x + 1 \pmod{x^8+1} .
	\end{equation*}
	Note that although $x^8+1$ is not irreducible (e.g. $x+1$ is a factor), $x^4 + x^3 + x^2 + x + 1$ is coprime to $x^8+1$ therefore it has an inverse$\mod{x^8+1}$. It follows that $\SubBytes$ is an invertible bijection.
	
	$\SubBytes$ is the only source of nonlinearity in AES thus it shall have excellent nonlinearity properties. This was the major theoretical weakness of DES, its predecessor. Practical linear attack against DES was introduced by Matsui \cite{matsui1994linear} and analyzed in detail and successfuly performed by Junod \cite{junod2001complexity}.
	
	\begin{note}
	\label{note:sboxtable}
		Since $\SubBytes$ has only $256$ possible inputs, it is usually precomputed and implemented as a lookup table.   %!% dát do appendixu tabulku?
	\end{note}

\subsubsection{ShiftRows}
	
	The $\ShiftRows$ operation is very simple -- it only cyclically shifts rows of the state array, $i$-th row is moved $i$ positions to the left, see example below. Note that rows are numbered from $0$. $\ShiftRows$ is obviously invertible.
	
	\begin{equation*}
	\arraycolsep=.5em\def\arraystretch{1.5}
		\begin{array}{|c|c|c|c|}
			\hline
			A_{0,0} & \cg{15} A_{0,1} & \cg{30} A_{0,2} & \cg{45} A_{0,3} \\
			\hline
			A_{1,0} & \cg{15} A_{1,1} & \cg{30} A_{1,2} & \cg{45} A_{1,3} \\
			\hline
			A_{2,0} & \cg{15} A_{2,1} & \cg{30} A_{2,2} & \cg{45} A_{2,3} \\
			\hline
			A_{3,0} & \cg{15} A_{3,1} & \cg{30} A_{3,2} & \cg{45} A_{3,3} \\
			\hline
		\end{array}
		\quad\xrightarrow{\ShiftRows}\quad
		\begin{array}{|c|c|c|c|}
			\hline
			A_{0,0} & \cg{15} A_{0,1} & \cg{30} A_{0,2} & \cg{45} A_{0,3} \\
			\hline
			\cg{15} A_{1,1} & \cg{30} A_{1,2} & \cg{45} A_{1,3} & A_{1,0} \\
			\hline
			\cg{30} A_{2,2} & \cg{45} A_{2,3} & A_{2,0} & \cg{15} A_{2,1} \\
			\hline
			\cg{45} A_{3,3} & A_{3,0} & \cg{15} A_{3,1} & \cg{30} A_{3,2} \\
			\hline
		\end{array}
	\end{equation*}
	
	The design motivation is to add inter-column diffusion. Note that $\ShiftRows$ does not diffuse bytes internally -- the following operation is here for this purpose.

\subsubsection{MixColumns}
	
	$\MixColumns$ operates on the state array in a column-wise manner. Let us pick a column from the state array and denote its elements by $A_0, A_1, A_2, A_3$, respectively. $\MixColumns$ considers it as a polynomial $a(z) = A_0 + A_1 z + A_2 z^2 + A_3 z^3$ -- an element of $F[z]$. Let $b(z)$ be output of $\MixColumns$, then
	\begin{equation}
	\label{eq:mixcolpoly}
		b(z) = c(z) \cdot a(z) \pmod{z^4+1}
	\end{equation}
	where $c(z) = \texttt{03}\cdot z^3 + \texttt{01}\cdot z^2 + \texttt{01}\cdot z^2 + \texttt{02}$. Note that here we used hexadecimal notation for elements of $F$ which could also be viewed as polynomials. To avoid confusion, we used $z$ instead of $x$.
	
	Note that $c(z)$ is coprime to $z^4+1$ therefore invertible$\mod{z^4+1}$. It follows that $\MixColumns$ is invertible.
	
	Unlike $\SubBytes$, $MixColumns$ operates on a substantially larger domain therefore cannot be directly implemented as a lookup table. Rather let us denote $b(z) = B_0 + B_1 z + B_2 z^2 + B_3 z^3$. Modular multiplication in Equation \ref{eq:mixcolpoly} which is performed in $F[z]$ can be rewritten into matrix multiplication back in $F$, indeed
	\begin{equation}
	\label{eq:mixcolmatr}
		\begin{pmatrix}
			B_0 \\ B_1 \\ B_2 \\ B_3
		\end{pmatrix}
		=
		\begin{pmatrix}
			\texttt{02} & \texttt{03} & \texttt{01} & \texttt{01} \\
			\texttt{01} & \texttt{02} & \texttt{03} & \texttt{01} \\
			\texttt{01} & \texttt{01} & \texttt{02} & \texttt{03} \\
			\texttt{03} & \texttt{01} & \texttt{01} & \texttt{02}
		\end{pmatrix}
		\begin{pmatrix}
			A_0 \\ A_1 \\ A_2 \\ A_3
		\end{pmatrix}.
	\end{equation}
	
	The goal of $\MixColumns$ is to introduce inter-byte diffusion within a column.

\subsubsection{AddRoundKey}
	
	Given a $128$-bit round key, the $\AddRoundKey$ operation simply bitwise XORs it on the state. The round key is derived from the encryption key using the $\KeySchedule$ routine which follows.
	
	$\AddRoundKey$ is the only step which utilizes keying material. It is its own inverse.

\subsubsection{KeySchedule}
	
	$\KeySchedule$ is a routine which expands $128$-bit encryption key into $11$-times longer $\ExpandedKey$. It will be indexed from $0$ by $128$-bit blocks.
	
	The most important thing for us is that $\ExpandedKey$ begins with plain encryption key. We omit further details here because $\KeySchedule$ is specified rather in algorithmic than algebraic manner, see \cite[pp.43-45]{daemen2013rijndael} for detail.


% ======================================================================
% ===   A E S   A L G O R I T H M   D E S C R I P T I O N            ===
% ======================================================================

\subsection{AES Algorithm Description}

AES, as already stated, consists of $10$ rounds of $4$ operations except for the last one where $\MixColumns$ is missing, see the following algorithm.

\begin{alg}
\label{alg:aes}
	Given a $128$-bit plaintext and key, return its AES encryption.
	\begin{algorithmic}[1]
		\Function{AES\_Encryption}{$Plaintext,Key$}
			\State $\ExpandedKey \gets \KeySchedule(Key)$
			\State $State \gets Plaintext$
			\State $\AddRoundKey(State, \ExpandedKey[0])$
			\For{$Round = 1 \to 9$}
				\State $\SubBytes(State)$
				\State $\ShiftRows(State)$
				\State $\MixColumns(State)$
				\State $\AddRoundKey(State, \ExpandedKey[Round])$
			\EndFor
			\State $\SubBytes(State)$
			\State $\ShiftRows(State)$
			\State $\AddRoundKey(State, \ExpandedKey[10])$
			\State\Return $State$
		\EndFunction
	\end{algorithmic}
\end{alg}

\begin{note}
	Since we have commented on every step about how it can be inverted, AES Encryption is also invertible by executing its steps inverted and in reverse order. This leads to AES Decryption algorithm.
\end{note}


% ======================================================================
% ===   A E S   I M P L E M E N T A T I O N   N O T E                ===
% ======================================================================

\subsection{AES Implementation Note}
\label{sec:aeslookup}

As already stated in Note \ref{note:sboxtable}, $\ShiftRows$ can be implemented as a lookup table for performance reasons. Such approach can be used for $\MixColumns$ as well. Let us denote
\begin{equation}
	\C =
	\begin{pmatrix}
		\texttt{02} & \texttt{03} & \texttt{01} & \texttt{01} \\
		\texttt{01} & \texttt{02} & \texttt{03} & \texttt{01} \\
		\texttt{01} & \texttt{01} & \texttt{02} & \texttt{03} \\
		\texttt{03} & \texttt{01} & \texttt{01} & \texttt{02}
	\end{pmatrix} .
\end{equation}
Since $\MixColumns$ is a linear operation, it follows from Equation \ref{eq:mixcolmatr}
\begin{equation}
	\begin{pmatrix}
			B_0 \\ B_1 \\ B_2 \\ B_3
		\end{pmatrix}
		=
		\C
		\begin{pmatrix}
			A_0 \\ 0 \\ 0 \\ 0
		\end{pmatrix}
		+
		\C
		\begin{pmatrix}
			0 \\ A_1 \\ 0 \\ 0
		\end{pmatrix}
		+
		\C
		\begin{pmatrix}
			0 \\ 0 \\ A_2 \\ 0
		\end{pmatrix}
		+
		\C
		\begin{pmatrix}
			0 \\ 0 \\ 0 \\ A_3
		\end{pmatrix}.
\end{equation}
Note that each addend only depends on a single byte i.e. there are just $256$ possible input values. Therefore $\MixColumns$ can be implemented as a fourtuple of lookup tables, the result is then a simple bitwise XOR of their $4$-byte output values.

This was just to give an idea about how AES operations can be turned into lookup tables. These can be further improved, see \cite[Chapter~4]{daemen2013rijndael} for detail.


\section{White-Box AES}
\label{sec:wbaes}

In this section we describe White-Box AES (WBAES) as introduced by Chow et al.\ \cite{chow2003aes} and revised by Muir \cite{muir2013tutorial} (which we highly recommend). Note that it has already been broken by Billet et al.\ \cite{billet2005cryptanalysis}, see Section \ref{sec:known}.   %!% ref to subsection, def linear?

The core idea of WBAES is based on usage of lookup tables as outlined in Section \ref{sec:aeslookup}, these are further improved in three ways:
\begin{enumerate}
	\item the implementation is turned into fully table representation i.e. the key is incorporated into tables as well,
	\item tables are wrapped around by appropriately self-vanishing random bijections, \label{item:wrap}
	\item the boundary is extended into the containing application. \label{item:boundary}
\end{enumerate}
Note that Chow et al. introduced also White-Box DES \cite{chow2003des} but it has been broken by Jacob et al.\ \cite{jacob2002attacking} much earlier than WBAES because it was lacking feature \ref{item:boundary} from the previous list.

\begin{notion}
\label{notion:table}
	Since WBAES will heavily use lookup tables, let us introduce a simplifying notion: a table will be told to have dimensions $m\times n$ if it has $m$-bit input values and $n$-bit output values. Note that the real size (i.e. memory consumption) of such table is $2^m\cdot n$ bits since you need $2^m$ entries of size $n$ to treat all possible $m$-bit input values.
\end{notion}


% ======================================================================
% ===   E N C O D I N G S                                            ===
% ======================================================================

\subsection{Encodings}

Let us first introduce some notion for feature \ref{item:wrap} from the previous list as defined by Chow et al.\ \cite{chow2003aes}.

\begin{defn}[Encoding]
	Let $X$ be a transformation from $m$ bits to $n$ bits. Choose an $m$-bit bijection $F$ and an $n$-bit bijection $G$. Call $X' = G \circ X \circ F^{-1}$ an {\em encoded} version of $X$. $F$ is an {\em input encoding} and $G$ is an {\em output encoding}.
\end{defn}

Note that, given a fixed key, AES itself is a $128$-bit bijection which can obviously never be implemented as a lookup table because its size would be $2^{128}\times 128$ bits. (On the other hand, this would be black-box equivalent -- you only know the input/output behavior.) Hence we will rather consider only a special kind of bijections -- we will use a series of smaller bijections and concatenate their outputs, a definition follows.

\begin{defn}[Concatenated encoding]
\label{def:concat}
	Let $F_i$ be $n_i$-bit bijections for $i=\atob{1}{k}$ and let $n = n_1 + \ldots + n_k$. The {\em function concatenation} $F_1 \| F_2 \| \ldots \| F_k$ is the $n$-bit bijection $F$ such that $F(b) = F_1(\atob{b_1}{b_{n_1}})\| F_2(\atob{b_{n_1+1}}{b_{n_1+n_2}})\|\ldots\| F_k(\atob{b_{n_1+\ldots+n_{k-1}+1}}{b_n})$ for any $n$-bit word $b=b_1b_2\ldots b_n$ where $\|$ denotes word concatenation. Clearly $F^{-1} = F_1^{-1}\| F_2^{-1}\| \ldots \|F_k^{-1}$.
\end{defn}

The following definition introduces the self-vanishing property of the bijections.

\begin{defn}[Networked Encoding]
\label{def:netw}
	A {\em networked encoding} for computing $Y\circ X$ (i.e. transformation $X$ followed by transformation $Y$) is an encoding of the form
	\[
		Y'\circ X' = (H\circ Y\circ G^{-1})\circ(G\circ X\circ F^{-1}) = H\circ(Y\circ X)\circ F^{-1}
	\]
	where $F$, $G$, $H$ are bijections of appropriate size, $F$ and $H$ are {\em external encodings} (input and output encoding, respectively) and $G$ is an {\em internal encoding}.
\end{defn}

Note that, given only tables for $Y'$ and $X'$, $G$ can be totally forgotten to be able to compute $Y\circ X$ provided that you know the input and output encoding. Such networked encoding can be obviously applied to much longer composition of transformations and you still only need to know the input and output encoding.   %!% někam popsat I/O encoding


% ======================================================================
% ===   R E O R D E R E D   A E S                                    ===
% ======================================================================

\subsection{Reorded AES}

In order to begin with fully-table representation, we need to reorder AES operations so that the resulting tables can be easily composed; let us first have a brief look at AES algorithm \ref{alg:aes}. Note that we can ``move'' the forcycle one row upwards with appropriate shift of indexes in $\ExpandedKey$. Then we switch $\ShiftRows$ and $\SubBytes$ without any side-effect since $\SubBytes$ operates byte-wise. Finally we switch $\ShiftRows$ and $\AddRoundKey$ while we moreover have to apply $\ShiftRows$ on $\ExpandedKey$ so that we get equivalent algorithm. The modified version looks as follows:
\begin{alg}
\label{alg:reordaes}
	~
	\begin{algorithmic}[1]
		\Function{Reordered\_AES\_Encryption}{$Plaintext,Key$}
			\State $\ExpandedKey \gets \KeySchedule(Key)$
			\State $State \gets Plaintext$
			\For{$Round = 0 \to 8$}
				\State $\ShiftRows(State)$
				\State $\AddRoundKey(State, \ShiftRows(\ExpandedKey[Round]))$
				\State $\SubBytes(State)$
				\State $\MixColumns(State)$
			\EndFor
			\State $\ShiftRows(State)$
			\State $\AddRoundKey(State, \ShiftRows(\ExpandedKey[9]))$
			\State $\SubBytes(State)$
			\State $\AddRoundKey(State, \ExpandedKey[10])$
			\State\Return $State$
		\EndFunction
	\end{algorithmic}
\end{alg}


% ======================================================================
% ===   F U L L Y - T A B L E   R E P R E S E N T A T I O N          ===
% ======================================================================

\subsection{Fully-Table Representation}

As outlined earlier, WBAES construction turns AES operations into table lookups only. Moreover we compose certain tables and wrap them all by internal encodings in the fashion of Definition \ref{def:netw}. WBAES generator thus inputs an AES key and a random seed and outputs key-dependent tables which serve as a WBAES instance.

\subsubsection{T-Boxes}
	
	First we compose $\AddRoundKey$ with the subsequent $\SubBytes$ step yielding an operation referred to as $\TBox$, also denoted by $T_{i,j}^r$ for $r$-th round and $(i,j)$-th position in the state array. We get
	\begin{equation}
		T_{i,j}^r(x) = S(x\xor k_{sr(i,j)}^r) , \qquad i,j=\atob{0}{3},\;r=\atob{0}{8}
	\end{equation}
	where $k_{sr(i,j)}^r$ stands for $\ExpandedKey$'s byte in $r$-th round at position $(i,j)$ shifted by $\ShiftRows$ (here denoted by $sr(i,j)$). The last round is treated separately as
	\begin{equation}
		T_{i,j}^{9}(x) = S(x\xor k_{sr(i,j)}^9) \xor k_{i,j}^{10} , \qquad i,j=\atob{0}{3} .
	\end{equation}
	
	Single $\TBox$ can be implemented as an $8\times 8$ lookup table\footnote{According to Notion \ref{notion:table}.} and there are $16\cdot 10 = 160$ of them.   %!% musim teda všude napsat kolik jich je

\subsubsection{MixColumns}
	
	$\MixColumns$ (here also denoted by $\MC$) can be turned into four $8\times 32$ lookup tables followed by three XOR operations as outlined in Section \ref{sec:aeslookup}. The difference from classical implementation is that the following XOR operations must be turned into lookup tables as well.
	
	Since the input size of such XOR operation would be $2\cdot 32$ bits, we first need to split inputs into smaller $4$-bit segments which can be treated separately. In the end we simply concatenate them obtaining the final result in the sense of Definition \ref{def:concat}. Single small XOR table is then $2\cdot 4\times 4$ and you need to perform $32/4\cdot 3=24$ such lookups per one column.

\subsubsection{ShiftRows}
	
	$\ShiftRows$ does not need table representation since it only moves bytes. We will see later how it can be implemented.


% ======================================================================
% ===   W B A E S   C O N S T R U C T I O N                          ===
% ======================================================================

\subsection{WBAES Construction}

So far we have trated all AES operations by lookup tables only. Now we will introduce a diffusion feature called {\em Mixing Bijection}, compose certain tables together and insert encodings which are actually a confusion feature.

\subsubsection{Mixing Bijection}
	
	Mixing bijection (denoted by $\MB$) is a random column-wise (i.e. $32$-bit) linear bijection inserted after $\MC$ and inverted in a separate step. Since it inputs the output of $\MC$, it can be composed together yielding $\MB\circ\MC$. Note that the table representation of both $\MB\circ\MC$ and $\MB^{-1}$ can be created in very the same manner as for $\MC$ itself.
	
	%!% 4x4 submatrices of full rank !!!

\subsubsection{Input Mixing Bijection}
	
	Input mixing bijection ($\IMB$) is a byte-wise variant of mixing bijection. It provides diffusion after inverting mixing bijection and before entering next $\TBox$. Since it is a linear mapping, it can be composed with inverted mixing bijection and since it is a byte-wise mapping it can be composed with $\TBox$ as well. The reason for this smaller diffusion element will be given later.

\subsubsection{Table Composition}
	
	Note that $\TBox$ is implemented as an $8\times 8$ table which is followed by an $8\times 32$ table of $\MB\circ\MC$. Therefore it is reasonable to compose those appropriately subsequent tables together -- we save both space and time, moreover we do not need to care about any byproducts.

\subsubsection{Inserting Encodings}
	
	Note that the (partial) $32$-bit outputs of linear operations are first split into $4$-bit words and then two independent $4$-bit words are XORed together yielding a new $4$-bit word. Thus we can only use $4$-bit internal encodings (denoted by $\Enc$) on both sides of each XOR. This is also the reason why input mixing bijection has been introduced -- it diffuses two $4$-bit blocks together.
	
	Let us demonstrate previous ideas in a single sketch of the flow during one inner round (i.e. not the first nor the last) showing all operations, table compositions, mixing bijections and encodings. Tables are given their names as introduced by Chow et al.\ \cite{chow2003aes}.
	\begin{align}
		\ldots&\rarr\underbracket{\Enc\rarr\IMB^{-1}\xrightarrow{\textnormal{plain}}\TBox\xrightarrow{\textnormal{plain}}\MB\circ\MC\rarr\Enc^{-1}}_{\textnormal{Type II}}\rarr\Bigl(\underbracket{\Enc\rarr\xor\rarr\Enc^{-1}}_{\textnormal{Type IV}}\Bigr)^2\rarr\nonumber\\[.5em]
		&\rarr\underbracket{\Enc\rarr\IMB\circ\MB^{-1}}_{\textnormal{Type III}}\rarr\Bigl(\underbracket{\Enc\rarr\xor\rarr\Enc^{-1}}_{\textnormal{Type IV}}\Bigr)^2\rarr\ldots
	\end{align}
	
	\ldots append Muir's figure depicting all operations together with their I/O sizes.
	%!% připojit Muirovo obrázek

\subsubsection{External Encodings}
	
	% Muir snad pochopil správně

%!% přidat minimálně zmínku o Karroumiho "vylepšení" včetně ekvivalence původnímu WBAES

\section{Known Attacks}
\label{sec:known}

As stated in Section \ref{sec:catmouse}, soon after introduction of a new white-box implementation of a cipher somebody usually comes up with an effective attack. % e.g. ... \cite introduced in ... was successfuly attacked by ... in ... \cite

This is the case of WBAES introduced by Chow et al. \cite{chow2003aes} in 2002 as well, it was broken by Billet et al. \cite{billet2005cryptanalysis} two years later. The attack is called {\em BGE attack} -- it is an abbreviation of authors' names.


% ======================================================================
% ===   B G E   A T T A C K                                          ===
% ======================================================================

\subsection{BGE Attack}

BGE attack is a pure algebraic attack therefore requires access to all tables implementing WBAES. % Note that we might need to reverse-engineer binary code first in order to access the tables.

Note that white-box version of DES by Chow et al. \cite{chow2003des} 
%~ DES že Boneh et al. zlomil kvůli slabším okrajům, it did not take long time and wb-aes was broken as well. ... je v tomhle odolnej, ale existuje na něj BGE útok, Karroumi "vylepšil" nevylepšil



		
		%%% KAPITOLA 2
		\chapter{Using Side-Channel Attack Tools}
\label{chap:attack}

In this chapter we first describe an attack exploiting information leaks assumed in gray-box attack context (see Section \ref{sec:bbtowb}). This attack is referred to as {\em Side-Channel Attack} (SCA) and we show two example SCA algorithms, both practically usable.

Next we present a groundbreaking approach of Bos et al. \cite{bos2015differential} utilizing side-channel attack tools in white-box attack context.

%!% kde je experiment design ve smyslu že mam sourcy a znam embedded key?

\section{Side-Channel Attack}
\label{sec:side}

%~ SCA, CPA, DCA

%~ Developed in ? by ?, correlation power analysis, differential computation analysis, \ldots


\section{Using SCA Tools in White-Box Cryptography}
\label{sec:scawbc}

%~ bude o útoku Bos et al. z teoretickýho pohledu, high level popis, vč. toho jak vzniknul
%~ trace acquisition, trace filtering, row & address filtering (theoreticaly describe the whole attack)

In the previous section we presented seemingly unrelated algorithms. Indeed, our aim is to attack white-box implementations where we can fully control the execution environment thus there is no need for physical measurements. The actual reason emerges from a groundbreaking approach introduced by Bos et al. \cite{bos2015differential} -- the use of SCA tools in white-box cryptography context.

The idea was inspired by Delerabl{\'e}e et al. \cite{delerablee2013white}. As it is usually required, a ``perfect'' white-box implementation should not provide effectively any other information about the secret key but only as much as an access to a black-box. In \cite{delerablee2013white} they observed that such implementation must be resistant to all existing and future side-channel attacks which led Bos et al. to the idea of exploring this consequence by attacking white-box implementations by SCA tools.

The obvious question is what kind of traces we should use instead of power traces -- we use {\em software traces}. In this context, software trace reffers to a record of memory addresses being accessed during program execution or their contents, sometimes also referred to as {\em memory trace}.

% RISCURE motivation:
%~ Back in November 2015, I had no results at all, therefore I waited until now, until I have some results. In my diploma thesis I am analyzing a novel attack published in a recent paper by Bos et al. \cite{bos2015differential}. Now I can successfuly reproduce the attack with own tools and identify the exact address locations where the keying material leaks.
%~ 
%~ Most interestingly, the attack breaks Klinec's implementation \cite{klinec2013white} of White-Box AES by Chow et al. \cite{chow2003aes}. To the best of my knowledge, there is no reasonable explanation why this is actually possible. The best known attack is BGE attack \cite{billet2005cryptanalysis} which is pure algebraic, on the other hand, Bos' attack utilizes side-channel tools to attack memory traces instead.
%~ 
%~ Even though it is not fully understood, the benefit of the attack is obvious -- it introduces a principially different approach to break white-box implementations. The attack may help to address weak points of an implementation and increase its resistancy to this kind of attack.
%~ 
%~ Since curious people tempt to break mysteries, I wanted to know why it works, too. I tried to set a watchpoint to the leaking address in a debugger and find the corresponding location in the source code. For some reason, standalone GDB did not catch the address previously caught by PIN and GDB connected to PIN debugger was not able to set any watchpoint at all. Hence I tried to print and attack all possible indexes used during encryption which should be later transformed to memory addresses, but surprisingly got absolutely nothing.

\subsection{Trace Acquisition}
\label{sec:tracq}

In order to acquire a memory trace we use Dynamic Binary Instrumentation (DBI) tools. These tools insert additional instructions to the original code of the program at run-time enabling one to debug or detect memory leaks. The most advanced DBI tools like Valgrind \cite{nethercote2007valgrind} or PIN \cite{luk2005pin} include a programmable interface where you can write tools of your will. Both tools are open-source.

We modified a tool for PIN by Teuwen \cite{teuwen2015movfuscator} for $4$ use cases. These acquire
\begin{enumerate}
	\item addresses of memory reads, \label{item:readaddr}
	\item addresses of memory writes,   %!% neni
	\item contents of memory reads, and   %!% neni
	\item contents of memory writes, \label{item:writecnt}
\end{enumerate}
respectively. Tool nr.\ \ref{item:writecnt} is obviously most useful for unprotected implementations since it can directly observe byproducts, tool nr.\ \ref{item:readaddr} will be surprisingly used as well.

Using chosen tool we acquire certain amount of memory traces for different, possibly random, plaintexts. We save these plaintexts along corresponding traces, of course, and proceed towards trace filtering.

\subsection{Trace Filtering}
\label{sec:filter}

Once we have acquired memory traces, we perform two kinds of filtering. We filter trace entries
\begin{itemize}
	\item by address and temporal range, and
	\item by constant value.
\end{itemize}

\subsubsection{Address and Temporal Filtering}

We have developed own tool displaying a memory trace. Given $n$, this tool splits the trace by address into $n$ ``compact'' segments and displays these parts separately. See Figure \ref{fig:memtrace} for an example memory trace. Horizontal axis represents address range, vertical axis represents the order of addresses from the top (temporal range).

\begin{figure}[h]
\begin{center}
	\includegraphics[width=0.9\textwidth]{./figures/memtrace/memtrace.png}
	\caption{Illustrative memory trace of a run of a simple AES implementation. Addresses of memory writes were acquired. Note $10$ AES rounds with $4$ inner steps.}
	\label{fig:memtrace}
\end{center}
\end{figure}

In this example trace we can see $10$ distinct AES rounds even with its $4$ inner steps. In general, if possible, we can estimate address and temporal range of the first round\footnote{Note that we are only interested in the first round, see Equation \ref{eq:byprod} and its context.} based on visual observation. Then we filter all traces according to this range. Note that this step brings the most substantial speedup of the attack since its speed depends on the size of traces.

\subsubsection{Constant Entries Filtering}

In the second stage we filter out such entries which are constant across all acquired traces because they do not carry any information. Indeed, both presented SCA algorithms exploit changes across traces.

Practically we first create a filter mask based on a small set of traces, $10$ shows to be sufficient. Such mask carries $0$ if all corresponding addresses in the small set of traces are equal, $1$ otherwise. This mask is then simply applied to filter all the remaining traces, possibly hundreds or thousands of them.

\subsection{Attacking Traces}
\label{sec:attack}



		
		%%% KAPITOLA 3
		\chapter{Practical Attack \& Results}
\label{chap:results}

First we describe what tools we needed to perform a practical attack, then we present our results and finally we discuss possible enhancements of the attack and how this attack can be useful for white-box cipher design.

\section{Attack Tools}

Since there is no open-source toolkit for the purpose of DCA described in Section \ref{sec:scawbc} in Chapter \ref{chap:attack}, we implemented it ourselves. Our toolkit consists of several tools and subtools, see the following list.
\begin{itemize}
	\item Acquire traces (see Section \ref{sec:tracq}):
	\begin{itemize}
		\item acquire single trace (PIN tools, modified versions of Teuwen's one \cite{teuwen2015movfuscator}),
		\begin{itemize}
			\item memory reads/writes and their addresses/contents,
		\end{itemize}
		\item manage trace acquisition (generate \& save respective plaintexts along).
	\end{itemize}
	\item Filter traces (see Section \ref{sec:filter}):
	\begin{itemize}
		\item filter constant entries \& save its mask $M_C$ (see Figure \ref{fig:constmask}),
		\item filter by address and temporal range,
		\begin{itemize}
			\item acquire another {\em full} trace $T$ (preferably in the same terminal session),
			\item filter it with mask $M_C$ yielding $T_C$,
			\item visualize $T_C$,
			\item estimate address \& temporal range,
			\item create another mask $M_R$ based on the range using $T_C$,
			\item filter all traces with $M_R$ (see Remark \ref{rem:rangemask}).
		\end{itemize}
	\end{itemize}
	\item Attack traces (see Section \ref{sec:attack}):
	\begin{itemize}
		\item run attack,
		\begin{itemize}
			\item CPA attack (Algorithm \ref{alg:cpa}), or
			\item bit-wise DPA attack (Algorithm \ref{alg:bitwisedpa}),
		\end{itemize}
		\item display results.
	\end{itemize}
\end{itemize}

First we used out toolkit to attack our straightforward implementation of AES which was used just as a proof-of-concept. Then we performed the most interesting attack Bos et al. \cite{bos2015differential} did -- we successfuly attacked Klinec's implementation \cite{klinec2013implementation} of WBAES by Chow et al. \cite{chow2003aes}.

%!% nápad: podle toho co vidim v memtrace, neni lepší útočit na 0x........??..; na 0x.........??. nebo dokonce na 0x........???? místo 0x..........??
% zkus vimdiff 0.txt 1.txt v attack/enhancements

\section{Results}

%~ co jsem použil za implmentace, kolik traců, jak dlouho to běželo, že jsem zjistil adresu kde to leakuje
%~ teoreticky by mělo bejt popsáno v předch. kapitole, tady už bych jenom zmínil co jsem použil a co vypadlo
%~ a proč to s tou javou nešlo ... (ale vyzkoušet)
%~ de facto takovej protokol
%~ Two (?) implementations will be considered for the attack:
%~ \begin{itemize}
	%~ \item our naive AES implementation (C++), and
	%~ \item Klinec's implementation of White-Box AES \cite{klinec2013white} (C++).
%~ \end{itemize}
%~ NaiveAES implementation was broken by both algorithms, CPA attack required $x$ traces, bitwise DPA attack required $y$ traces.
%~ White-Box-AES implementation was only broken by bitwise DPA attack and many interesting things occured.
% provést útok s několikrát nagenerovanejma tabulkama (v manuálu pak popsat co je pro to potřeba udělat)

\section{Enhancements of the Attack}

%~ zmínka o inverzi, že zas nic ...


%?% vlastní kapitola?
\section{Use in White-Box Crypto Design}

%~ pokud bych našel v kódu kde to leakuje, uměl to vyčíst z indexů popř. nák přímo z tabulek -> 4 Analysis

% chtěl jsem pokračovat ve smyslu:
% 	najít debuggerem místo v kódu, kde to leakuje, a zjistit tak kterej index to je, případně se to naučit počítat
% ale nedařilo se, přestože jsem vyzkoušel několik způsobů:
% 	pustil jsem to v gdb, to ale leakující adresu vůbec nezachytilo (! záleží na input plaintextu, ohlídal jsem si to)
% 	zkusil jsem gdb připojit na pin, ale to nešlo vůbec, hned hlásí že Could not insert hardware breakpoints: You may have requested too many hardware breakpoints/watchpoints.
% 	další mě nenapadá ... (zkusit valgrind a porovnat tracy?)
% alternatively, přímo vypsat všechny možný indexy použitý v tabulkách a zkusit útočit na tohle místo memtrace
% 	z toho poskládat přímej útok

% proposed countermeasures against DCA?
%~ To the best of my knowledge, there is no reasonable explanation why this is actually possible. The best known attack is BGE attack \cite{billet2005cryptanalysis} which is pure algebraic, on the other hand, Bos' attack utilizes side-channel tools to attack memory traces instead.
%~ Even though it is not fully understood, the benefit of the attack is obvious -- it introduces a principially different approach to break white-box implementations. The attack may help to address weak points of an implementation and increase its resistancy to this kind of attack.
%~ I tried to set a watchpoint to the leaking address in a debugger and find the corresponding location in the source code. For some reason, standalone GDB did not catch the address previously caught by PIN and GDB connected to PIN debugger was not able to set any watchpoint at all. Hence I tried to print and attack all possible indexes used during encryption which should be later transformed to memory addresses, but surprisingly got absolutely nothing.

		
		%%% KAPITOLA 4
		%~ \chapter{Analysis of the Attack against WBAES}
\label{chap:analysis}

In this chapter we present some remarks of the attack introduced in Section \ref{sec:unify} against Klinec's implementation \cite{klinec2013implementation} of WBAES by Chow et al.\ \cite{chow2002aes}, denoted as {\tt KlinecWBAES}. Based on our observations, we suggest blind attack (i.e.\ an attack without knowledge of the key) and mention some consequences in white-box design. % Then we outline why this attack is useful during white-box cipher design.
Finally we outline, justify and confirm an explanation of the attack.


\section{Blind Attack}
\label{sec:blindattack}

In a real-world scenario, the key is typically unknown. Therefore we do not have any information whether our best candidate is correct or not. We need some rule to recognize a fairly good candidate, which we can then test for correctness by running AES encryption ourselves.

The most straightforward rule seems to be thresholding the gap of the best candidate, but it turns out that it is not that easy, see the following remark.

\begin{remark}
\label{rem:false}
	Especially note the attack against the \nth{1} byte using {\tt 3d} as a target and its last but one entry in Table~\ref{tab:lintargets} -- here an incorrect candidate has a gap of almost $26\%$. It is actually the largest gap of an incorrect candidate among all targets and all bytes using this particular instance of WBAES tables. (The largest gap of an incorrect candidate ever seen was almost $35\%$!)
	
	In our results, the target bits keep their original order (but reverse), therefore we can deduce the vector $B$ generating this particular target in terms of $T_B$ as introduced in Equation~\ref{eq:tba}. Here it is the second row of matrix of multiplication by $\texttt{3d}\pmod{x^8+1}$, see Example~\ref{ex:shiftmatrix}. Hence the vector is $B = 01011110$.
\end{remark}

In order to perform a blind attack, we studied our results deeper -- we were particularly interested in success rates of individual targets. The following section gives some of our remarks, next we suggest the blind attack itself in Section~\ref{sec:subblindattack}


% ==============================================================================
% ===   R E M A R K S                                                        ===
% ==============================================================================

\subsection{Remarks}
\label{sec:remarks}

%~ We decided to create another {\tt KlinecWBAES} table instantiations in order to have larger data set. Note that for each instance we have $16$ independent tables -- one for each byte since %?% neni to podle mě uplně nezávislý ptž MC je před MB tže dycky aspon 4 byty jsou provázaný
We obviously cannot make conclusions based on a single {\tt KlinecWBAES} table instantiation only. In order to avoid specific behavior of a single instance, we created another $7$ instances. Then we acquired $1024$ traces and ran the attack against each of $16$ key bytes using all $255$ targets, for each instance. Altogether we ran $16\cdot255\cdot8=32\,640$ attacks, which, together with trace acquisition and filtering, took us several hours on a regular hardware.

We processed the results and displayed several statistics, here we present some of them. Note that we only considered strong candidates as introduced in Note~\ref{note:strong} (candidates with a gap greater than $10\%$).

The good news is that the overal success rate was $25\%$ (cf.\ $27\%$ and $31\%$ for our previous single-instance attack using the original SBox and Rijndael inverse, respectively), which gives us a hope that the new targets are successful, too. But our goal is to estimate rather individual success rate of each of our $255$ targets -- here the good news is that each target was successful at least once out of $16\cdot8=128$ chances.

Let us give the remarks now -- namely we present success rates of our targets, behavior of incorrect candidates, what happens if we use less than $1024$ traces, and finally we show a very peculiar result regarding leaking position within our traces.

\subsubsection{Success Rate of Individual Targets}
	
	%!% přidat footnote, že ještě nevíme jak jsou nezávislý
	Remind that we have $255$ targets and only $8$ independent instances of tables which is quite a small data set for statistical purposes. For this reason, we rather group targets by different criteria and observe if they appear to be uniformly distributed. Note that we could possibly only refute uniformity this way, but it can still serve as reasonable justification.
	
	\paragraph{Group by corresponding $p$.}
	
	%~ We studied whether there is some influence of target invertibility (non-invertible targets were introduced in Section~\ref{sec:noninv}).
	
	Here we make use of our former approach, and group the targets by corresponding $p$; see the transfer in Section~\ref{sec:unify}. We put the results in a histogram, see Figure~\ref{fig:leaktargethist}. Here we give average percentual success rate together with its standard deviation (among $8$ measurements), for each group of targets. Invertible $p$'s are shown in green, non-invertible ones are orange.
	
	\begin{figure}[h]
	\begin{center}
		\includegraphics{figures/leak_target/leak_target.pdf}
		\caption{Average success rate and its standard deviation for targets grouped by corresponding $p$. Green represents invertible $p$'s, orange noninvertible ones.}
		\label{fig:leaktargethist}
	\end{center}
	\end{figure}
	
	Note that we did not give $y$-axis scale, since the only purpose is to distinguish uniform distribution, which appears to achieved be here.
	
	\paragraph{Group by fixed $4$ bits of $B$.}
	
	Another way of grouping we performed was grouping targets by fixing $4$ out of $8$ bits of their corresponding vector $B$. We fixed the first and the last $4$ bits and grouped them; see results in Figure~\ref{fig:leaktargetotherhist}, where these ways of grouping are given in green and orange, respectively. Note that group index can be computed as binary AND of mask {\tt 0xf0} and {\tt 0x0f} with vector $B$, respectively.
	
	\begin{figure}[h]
	\begin{center}
		\includegraphics{figures/leak_target_other/leak_0x0f_0xf0.pdf}
		\caption{Average success rate and its standard deviation for targets grouped by fixed $4$ bits of corresponding vector $B$. Green represents fixed first $4$ bits, orange represents fixed last $4$ bits.}
		\label{fig:leaktargetotherhist}
	\end{center}
	\end{figure}
	
	\begin{remark}
	\label{rem:uniform}
		According to our previous observations, there does not seem to be a preferred target or a group of targets, therefore we will assume that success rate is uniform across all targets.
	\end{remark}

\subsubsection{Incorrect Candidates}
	
	There is an inconvenience, which emerges once we do not know the key -- we do not recognize the correct candidate anymore, see Remark~\ref{rem:false}. Therefore we will first observe behavior of incorrect ones in our self-controlled experiment where we know the key.
	
	\paragraph{Gaps of incorrect candidates.}
		
		Previously we defined a strong candidate as a candidate with a gap greater than $10\%$. Most strong candidates were correct, too, but there were a couple of exceptions, which we will refer to as {\em false positives}. In our overall results, $25\%$ of candidates were strong and correct, on the other hand, there were $11\%$ of false positives. On average, false positives had a gap of $14\%$ (cf.\ $40\%$ for correct ones) and the highest gap ever seen was $35\%$ (cf.\ $76\%$ for correct ones).
		
		A criterion based on thresholding might be on the one hand too strict, on the other hand, it might fail. Instead of looking for some well-balanced threshold, let us have a look at a much better observation about false positives, which will help us to distinguish them easily from the correct ones.
		% 1024:
		% strong & correct:
		% 	[1019, 1000, 978, 1008, 986, 1006, 1019, 999] => 1001.875 ... 24.6%
		% 	[39.46, 39.68, 39.5, 40.28, 40.2, 39.85, 39.86, 39.68] => 39.81375
		% strong & incorrect:
		% 	[442, 440, 425, 415, 433, 463, 413, 468] => 437.375 ... 10.7%
		% 	[14.05, 14.06, 14.16, 13.95, 14.26, 13.83, 14.0, 13.97] => 14.035
	
	\paragraph{Number of repetitions of false positives.}
		
		The good news is that the same false positive does not appear to repeat across targets very often. We observed the number of most repeating false positives across our $255$ targets and got the following results: on average, the number was $1.75$, the global maximum was $3$.
		
		Therefore having a candidate which repeats several times, we most likely have the correct candidate.
		% 	[1.88, 1.75, 1.81, 1.81, 1.62, 1.88, 1.62, 1.62]

\subsubsection{Using Less Traces}
	
	So far, we always used $1024$ traces in our attacks, let us have a look at results of the attack with much less traces. Note that Bos et al.\ \cite{bos2015differential} used $2000$ and $500$ traces for the attack targetting the original SBox and Rijndael inverse, respectively.
	
	Namely, we used $128$, $256$, $384$ and $512$ traces. For each number of traces, we attacked all of our $8$ instances of {\tt KlinecWBAES}, and observed ratios of strong candidates (both correct and incorrect) and their average gap; see results in Table~\ref{tab:ntraces}. Note that number of repeating false positives remains very low -- up to three.
	
	\begin{table}[H]
		\begin{center}
		\begin{tabular}{| c | c | c | c | c | c | c | c |}
			\hline
			Traces &    $128$ &    $256$ &    $384$ &    $512$ &    $1024$ \\
			\hline
			\hline
			Correct candidates
			       &  $9.8\%$ &   $17\%$ &   $20\%$ &   $21\%$ &    $25\%$ \\
			\hline
			Average gap of correct candidates
			       &   $24\%$ &   $32\%$ &   $35\%$ &   $37\%$ &    $40\%$ \\
			\hline
			\hline
			False positives
			       &  $7.7\%$ &   $11\%$ &   $12\%$ &   $12\%$ &    $11\%$ \\
			\hline
			Average gap of false positives
			       &   $13\%$ &   $14\%$ &   $14\%$ &   $14\%$ &    $14\%$ \\
			\hline
			\hline
			Reduced cost\tablefootnote{Reduced cost is to be introduced later.}
			       & $5\,400$ & $4\,700$ & $5\,500$ & $6\,400$ & $10\,500$ \\
			\hline
		\end{tabular}
		\end{center}
	\caption{Ratios of correct and incorrect candidates and their average gaps, for different numbers of traces.}
	\label{tab:ntraces}
	\end{table}
	
	% 128:
	% strong & correct:
	% 	[388, 420, 391, 411, 402, 393, 412, 397] => 401.750 ... 9.8%
	% 	[24.23, 24.39, 24.26, 23.65, 24.5, 24.01, 23.87, 23.65] => 24.07
	% strong & incorrect:
	% 	[262, 325, 351, 322, 309, 342, 285, 315] => 313.875 ... 7.7%
	% 	[13.34, 13.34, 13.13, 13.12, 12.9, 13.22, 13.26, 13.33] => 13.205
	
	% 256:
	% strong & correct:
	% 	[700, 670, 658, 722, 700, 709, 707, 683] => 693.625 ... 17.0%
	% 	[31.74, 32.49, 32.41, 31.83, 32.29, 31.31, 31.91, 32.11] => 32.01125
	% strong & incorrect:
	% 	[431, 472, 423, 441, 447, 451, 424, 482] => 446.375 ... 10.9%
	% 	[13.82, 13.85, 13.69, 14.1, 14.01, 14.1, 13.95, 14.36] => 13.985
	
	% 384:
	% strong & correct:
	% 	[814, 791, 784, 825, 790, 819, 812, 818] => 806.625 ... 19.8%
	% 	[35.84, 35.59, 35.14, 35.26, 36.19, 35.26, 35.77, 34.61] => 35.4575
	% strong & incorrect:
	% 	[491, 488, 484, 504, 487, 464, 509, 530] => 494.625 ... 12.1%
	% 	[14.17, 14.37, 13.97, 14.17, 14.14, 14.07, 14.23, 14.3] => 14.1775
	
	% 512:
	% strong & correct:
	% 	[879, 851, 855, 887, 848, 879, 881, 870] => 868.75 ... 21.3%
	% 	[37.25, 37.47, 36.58, 37.66, 38.04, 37.03, 37.42, 36.93] => 37.2975
	% strong & incorrect:
	% 	[471, 480, 473, 478, 472, 517, 432, 474] => 474.625 ... 11.6%
	% 	[14.02, 14.09, 13.96, 13.97, 14.23, 13.71, 14.29, 13.97] => 14.03

\subsubsection{Within-Trace Position of Leaks}
	
	Remind that our trace is a bit-wise serialization of least significant bytes of memory addresses, hence we can identify which bit within that byte was leaking -- simply by taking the leaking position within trace modulo $8$.
	
	\begin{note}
	\label{note:leakbit}
		This moduled position will be referred to as {\em leaking bit} (note the difference from target bit).
	\end{note}
	
	We noticed soon that leaking bits are not very well balanced as one would expect to, therefore we put this data into a histogram; see Figure~\ref{fig:leakbitall}. The histogram shows, for each leaking bit, the average number of strong correct candidates together with its standard deviation, which is surprisingly very low. Note that the histogram does not provide $y$-axis scale for similar reasons as for previous histograms.
	
	\begin{figure}[h]
	\begin{center}
		\includegraphics{figures/leak_bit/leak_bit.pdf}
		\caption{Average number of leaks and its standard deviation at each leaking bit.}
		\label{fig:leakbitall}
	\end{center}
	\end{figure}
	
	The \nth{0} and \nth{1} bits leaked slightly more, the \nth{2} bit slightly less than average, and the \nth{3} bit actually leaked only twice, while the overal average of the remaining bits was almost $160$! On the other hand, all of the remaining bits leaked fairly similarly. We do not have any explanation for this behavior, we only present it as our observation.


% ==============================================================================
% ===   B L I N D   A T T A C K   S U G G E S T I O N                        ===
% ==============================================================================

\subsection{Blind Attack Suggestion}
\label{sec:subblindattack}

\begin{note}
	This section only applies previous observations on a heuristic basis, there is no guarantee that our approach is the best.
\end{note}

In general, we suggest to use rather less traces, and repeat the attack with several targets until the sum of strong gaps of any candidate exceeds given bound. Let us choose reasonable values for both the number of traces and the bound.
\begin{description}
	\item[Number of traces.]
		In order to deduce a reasonable number of traces from results in Table~\ref{tab:ntraces}, let us introduce {\em reduced cost of gap}, defined as
		\begin{equation*}
			C(n, s, g) = \frac{n}{s\cdot g} ,
		\end{equation*}
		where $n$ stands for the number of traces, $s$ for average success rate and $g$ for average gap of a strong candidate.
		
		Note that this quantity corresponds with average time of the attack. Indeed, the more traces, the longer time; the better success rate or the bigger gap, the shorter time. Hence we can use this quantity to compare expected time of the attack using different number of traces.
		
		According to results in Table~\ref{tab:ntraces}, the best value of the reduced cost of gap was achieved for $256$ traces, therefore we suggest to use $256$ traces.
	\item[Bound.]
		According to our previous observation, false positives do not appear to repeat very much -- there were only a couple of repetitions for the most successful ones --, for less traces neither. Therefore we suggest to use $75\%$ as a cummulative bound, since it requires some $6$ hits for a fairly successful false positive, which appears to be hardly achievable.
\end{description}

\begin{remark}
\label{rem:attimpr}
	The attack can be further improved -- we can make the bound adaptive: we begin with a lower bound, obtain some (full) key candidate and verify it. Until the candidate is correct, we keep increasing the bound and merge our previous and new results. However, we did not implement this enhanced variant.
	%?% vyzkoušet jestli tak zlomim všechny instance
\end{remark}

% psal já / Teuwen:
%~ > I only wonder about the reasoning why Karroumi is more than Chow since
%~ > it seems to have been shown to be equal (based on what I wrote in my
%~ > previous email).
%~ 
%~ Well I've no problem to break completely Chow with standard DCA so there
%~ is something a bit more in Karroumi. Obviously not enough to make it
%~ robust enough...


\section{Use in White-Box Cryptography Design}
\label{sec:useindesign}

So far we were only interested in results of the attack. In this section, we will study deeper the attack itself. First we describe our effort trying to identify which intermediate product causes the leakage, then we outline some countermeasures against this attack. Both might be useful in white-box cryptography design.


% ==============================================================================
% ===   L E A K I N G   I N T E R M E D I A T E   P R O D U C T              ===
% ==============================================================================

\subsection{Leaking Intermediate Product}

Note that the attack allows us to reveal concrete addresses which contributed to the key recovery. Our intention was to find the corresponding piece of source code and later possibly also the intermediate product which is later transformed into memory address -- it would most likely be an array index. We tried several approaches how to find that place in source code, however, none of the debugging approaches worked.

\begin{description}
	\item[Watchpoint in GDB.]
		We acquired one full trace (as described in Section \ref{sec:filter} under the caption ``\AddrTempFilter'') for some plaintext used during the acquisition phase. We checked that this trace perfectly matches the corresponding trace from the acquisition phase.
		
		We identified the leaking address and set up a watchpoint in GDB running in the same terminal session. However, GDB did not catch the address, probably because PIN uses different virtual environment.
	\item[GDB connected to PIN's debugging interface.]
		We set up a connection between PIN application debugging interface and GDB as described in PIN manual \cite{pin214manual}, but could not catch {\em any} address. GDB was complaining that it cannot insert hardware breakpoint, on the other hand, software breakpoints did not work either.
	\item[Outputting all array indexes.]
		Note that array indexes are very likely the values which are later transformed into leaking memory addresses. Hence it should be possible to output and attack all indexes used within the encryption procedure. However we did not succeed with the first attempt -- it required certain reasoning and finding the specific intermediate result to output; both will be given in Section \ref{sec:attempt}.
\end{description}

Each of the debugging approaches seemed promising to bring valuable information about where the key leaks. The leakage position could be later used to identify a specific table access which causes the leakage, hence it could possibly help us to clarify what is the actual intermediate product which enables the attack. However, we succeeded even without this information.

%!% překopat
Even though it is not fully understood, the benefit of this attack is obvious -- it introduces a principially different approach to break white-box cryptography implementations. This attack may help to address weaknesses of a future white-box cryptography implementation and possibly increase its resistancy to this kind of attack.


% ==============================================================================
% ===   C O U N T E R M E A S U R E S   A G A I N S T   D C A                ===
% ==============================================================================

\subsection{Countermeasures against DCA}

Since DCA is based on an algorithm which was originally developed for physical measurements of certain hardware emissions, we can look for some inspiration back into hardware model. There are several countermeasures, see the following list for some of them (basic ones are introduced in \cite{chari1999towards,goubin1999des}).
\begin{description}
	\item[Masking with random values.] Cryptographic hardware can run some unpretendable source of random data and thus mess up values in traces (i.e.\ adding a strong noise). However, there is no equivalent of noise in white-box context since everything is is fully controlled by the attacker.
	\item[Reordering instructions.] Remind that our algorithms required aligned traces, therefore it would be fatal if the leaking position were at several different places. However, this could be possibly overcome in white-box context -- we could achieve this by reverse engineering or, much easier, by aligning the traces with respect to instruction address trace (as outlined by Bos et al.).
	\item[Adding random delays.] Note that random delays have very similar effect as the previous countermeasures.
\end{description}
In general, we cannot rely on any source of {\em dynamic} random data (generated during program execution). All we need to rely on is {\em static} random data (generated during instantiation). On the other hand, dynamic randomness can be used in a white-box implementation -- just as another level of ``obfuscation''.

Bos et al.\ further propose to use some ideas from {\em threshold implementations} \cite{nikova2006threshold} or use of external encodings -- here they emphasize that more research is needed.


\section{Explanation Attempt}
\label{sec:attempt}

First of all, let us remind the flow of information through WBAES tables (originally in Equation \ref{eq:wbaesflow}), especially its first part.
\begin{equation}
\label{eq:wbaesfirst}
	\ldots \rarr \underbracket{\Enc \rarr \IMB^{-1} \xrightarrow{\textnormal{plain}} \TBox \xrightarrow{\textnormal{plain}} \MB \circ \MC \rarr \Enc^{-1}}_{\textnormal{in table}} \rarr \ldots
\end{equation}
Note that the last plain AES state (denoted with ``plain'' in the previous equation) is actually an $8$-bit output of the first SBox i.e.\ the very original target. Let us see what happens next: a composed linear mapping $\MB\circ\MC$ is applied, resulting in a $32$-bit output, and then each $4$-bit nibble is passed through a $4$-bit random bijection $\Enc^{-1}$. 

Note that $\Enc$ is only a $4$-bit bijection 

%!% attempt to explain: výstup SBoxu je problitej náhodnym lin. zobrazenim tj. prvky výstupu jsou zase jenom skal. součinem s nějakým náhodnym nenulovym vektorem ... po čtveřicích jsou takový bity problitý náhodnou bijekcí (právě že jenom po 4!)
% nikdo přece netestuje nelinearitu/nekorelovanost techdle bijekcí ... neni to třeba tak že každá 3. koreluje na nákym výstupnim bitu s nákym vstupnim bitem? nebo klíďo jejich lin. kombinací?
% ale mam pocit že to jde ověřit právě výpisem mezivýsledků a útokem na ně, co si pomatuju tak to nešlo ...

% ze dřívějška:
% neni třeba problém v tom že to jsou concatenated encodings? i.e. ze 2x4 bitů postavený 8 bitový bijekce
% jesli takovejch už třeba neni moc málo nebo jestli náhodou nedávaj docela často nákou korelaci
% třeba zkusit dokázat, že random bijection ani náhodou nezanechá korelaci v nákym bitu, mělo by tam bejt negligible
% asi se trochu rozepsat a odůvodnit ...
% třeba to, jak jsou ty bijekce generovaný, že to je z nákýho pejpru, kterej něco ukazuje


   % kéž by ...
		
		%%% KAPITOLA 5
		\chapter{Future Work}
\label{chap:future}

\section{Another Interesting AES White-Box Implementation}
\label{sec:magicimpl}

Bos et al.\ \cite{bos2015differential} mention another interesting AES white-box implementation in their paper. This particular implementation by Eloi Vanderbéken was first published only as a Windows binary executable challenge at NoSuchCon 2013 conference\footnote{\link{http://www.nosuchcon.org/2013/}}, where nobody was able to recover the key. Note that this implementation uses an externally-encoded variant of AES (as outlined in Section \ref{sec:constrwbaes}, in particular in Equation \ref{eq:extenc}), hence it is {\em not compliant with AES anymore}. It follows that any SCA-like attack is impossible, since the internal AES handles encoded inputs, where the encoding is unknown.

Soon after the conference, the author released source codes\footnote{\link{http://pastebin.com/MvXpGZts}}, which allows one to avoid reverse engineering of the binary. Moreover, it allows one to modify the code to run natively on Linux. Philippe Teuwen, a co-author of Bos et al.\ \cite{bos2015differential}, studied this implementation deeper and finally was able to recover the key \cite{teuwen2015nscwriteups}.

In his write-up, Teuwen also describes the structure, which appears to avoid usage of $4$-bit random bijections (in WBAES by Chow et al.\ \cite{chow2002aes}) and uses $8$-bit ones instead. Note that even without any non-linearity check, it is far less likely that any bit of output of a random $8$-bit bijection would provide a reasonable level of linearity compared to a random $4$-bit bijection, as outlined in Section \ref{sec:attempt}.

However, we did not either modify the source code to make it AES-compliant, nor did run the attack, hence it remains as a possible topic for further research.

%?%
% \section{What next?}
		
		%%% ZÁVĚR
		% cleardoublepage and phantomsection force that link jumps ABOVE heading, normally it would jump UNDER heading
\cleardoublepage\phantomsection

\addcontentsline{toc}{chapter}{Conclusion}
\chapter*{Conclusion}
\markboth{\MakeUppercase{Conclusion}}{}
\label{chap:concl}

%~ \item[Chapter~\ref{chap:intro}: \nameref{chap:intro}.] ~ \\
	%~ First of all, we give an introduction to cryptography and white-box cryptography. Then we present the structure of Advanced Encryption Standard (AES) -- a cipher which we will be particularly interested in during the whole thesis. Finally, we conclude this chapter with a description of a particular design of a white-box variant of AES.
After a general introduction to cryptography, we presented three attack contexts: black-box, gray-box and white-box, with emphasis on the last-named. Then we described another pillar of this thesis -- the Advanced Encryption Standard (AES) and its white-box variant by Chow et al.\ \cite{chow2002aes} (WBAES).

%~ \item[Chapter~\ref{chap:attack}: \nameref{chap:attack}.] ~ \\
	%~ Next we present an attack technique which originally targets hardware leakage. However, it can be advantageously used in software scenario as well -- we describe a recent attack based on this idea.
Next we presented side-channel attack which seems to be unrelated to white-box attack context. However, we described a relationship as originally outlined by Bos et al.\ \cite{bos2015differential}, and provided a detailed description of the attack. Here we especially ephasized one practical consequence -- no reverse engineering effort is needed for this attack.

%~ \item[Chapter~\ref{chap:results}: \nameref{chap:results}.] ~ \\
	%~ In the third chapter, we first reproduce previous results of the attack, then we suggest our improvements and run the attack again. Finally we introduce an approach that unifies our contributions.
In Chapter~\ref{chap:results}, we first reproduced results of the attack by Bos et al.\ against Klinec's implementation \cite{klinec2013implementation} of WBAES. Then we extended their attack from $16$ to $255$ targets, and finally we introduced a simple unifying approach for our contributions.

%~ \item[Chapter~\ref{chap:analysis}: \nameref{chap:analysis}.] ~ \\
	%~ In the fourth chapter, we present some remarks regarding the attack, then we exploit them to suggest a blind attack (i.e., an attack without knowing the actual key). We also mention some consequences in white-box implementation design. To conclude this chapter, we outline, justify and confirm our explanation of the attack.
In the fourth chapter, we analyzed the attack from three perspectives. First, we suggested a blind attack (i.e., an attack without knowing the actual key) based on our observations. Second, we commented on usage of the attack in white-box cryptography design and presented possible practical countermeasures which, however, did not seem to avoid this attack. Third, we identified and confirmed which intermediate result causes the leakage. We also outlined which single element protects WBAES from our enhanced attack and why its choice is inappropriate.

%~ \item[Chapter~\ref{chap:future}: \nameref{chap:future}.] ~ \\
	%~ Finally, we summarize possible topics of further research.
We concluded this thesis in Chapter \ref{chap:future}, where we summarized possible topics of further research.

		
		%%% LITERATURA
		\cleardoublepage\phantomsection   % cleardoublepage a phantomsection udělají že odkaz z obsahu vede NAD nadpis, jinak vede POD
\addcontentsline{toc}{chapter}{References}

\bibliographystyle{plain}
\bibliography{biblio}

		
		%%% APPENDIX
		\setcounter{tocdepth}{0}
\begin{appendices}

\chapter{Contents of Attached CD}
\pagenumbering{Roman}
\label{app:cd}
	
	\begin{description}
		\item[\tt diploma-thesis.pdf] ~ \\ diploma thesis in Portable Document Format,
		\item[\tt attackTools/] ~ \\ folder with our attack tools,
		\begin{description}
			\item[\tt attack/] ~ \\ our attack toolkit written in Ruby, see Section~\ref{sec:tools},
			\item[\tt BitwiseDPA/] ~ \\ C++ sources of customized bitwise DPA attack, see Algorithm~\ref{alg:bitwisedpa},
			\item[\tt MemoryTracingTools/] ~ \\ C++ sources of memory tracing tools, see Section~\ref{sec:tracq},
		\end{description}
		\item[\tt implementations/] ~ \\ folder with (white-box) AES implementations and tables,
		\begin{description}
			\item[\tt klinecWBAES/] ~ \\ C++ sources of customized {\tt KlinecWBAES}, see Section~\ref{sec:results},
			\item[\tt naiveAES/] ~ \\ C++ sources of {\tt naiveAES}, see Section~\ref{sec:results},
			\item[\tt NoSuchCon2013Challenge/] ~ \\ C source of an interesting AES white-box implementation, see Section~\ref{sec:magicimpl},
			\item[\tt WBAESTables/] ~ \\ our particular $8$ instances of {\tt KlinecWBAES} tables, see Section~\ref{sec:remarks}.
		\end{description}
	\end{description}


\end{appendices}
		
		
%%%%%%%%%%%%%%%%%%%%%%%%%%%%%%%%%%%%%%%%%%%%%%%%%%%%%%%%%%%%%%%%%%%%%%%%
%%%   K O N E C   P R Á C E
%%%%%%%%%%%%%%%%%%%%%%%%%%%%%%%%%%%%%%%%%%%%%%%%%%%%%%%%%%%%%%%%%%%%%%%%


\end{document}